%% ============================================================================
%% Chapter 3: The REST Architectural Style & Resource Design
%% ============================================================================
\chapter{The REST Architectural Style \& Resource Design}
\label{ch:rest_style}

\section*{Executive Summary}

Representational State Transfer (REST) is the most influential --- and the most
misunderstood --- architectural style in modern software engineering. Coined by Roy
Fielding in his 2000 doctoral dissertation \cite{fielding2000rest}, REST is not a
protocol, not a specification, and emphatically not ``JSON over HTTP.'' It is a
\emph{style}: a named set of architectural constraints that, when applied together,
induce a specific, desirable set of system properties --- scalability, reliability,
visibility, portability, and independent evolvability. An HTTP API earns the label
``RESTful'' only to the degree that it honors those constraints; an HTTP API that
ignores them is merely an HTTP API, which may still be a perfectly defensible
engineering choice --- provided the choice is made deliberately rather than by
accident of vocabulary.

This chapter formalizes Fielding's six constraints (client--server, statelessness,
cacheability, uniform interface, layered system, and the optional code-on-demand),
stating each precisely, enumerating the properties it induces, and identifying what
each forbids. We then confront the industry's chronic mislabeling problem through the
Richardson Maturity Model \cite{richardson2013restful}, a diagnostic ladder from the
``swamp of POX'' (Plain Old XML) through resources and HTTP verbs up to hypermedia
controls. The middle of the chapter is devoted to the discipline of resource design:
identifying resources as nouns, distinguishing collections, documents, stores, and
controllers, and engineering URI grammars that are hierarchical, predictable, and
opaque to clients. We give the uniform interface its full semantic treatment ---
method safety and idempotency, precise status codes for every CRUD lifecycle
transition, PUT versus PATCH (including JSON Merge Patch and JSON Patch), and content
negotiation --- before turning to HATEOAS and the practical economics of hypermedia
media types (HAL, JSON:API, Collection+JSON). Representation design closes the
theory: embedding versus linking, partial representations, and the anti-corruption
discipline that keeps your API from becoming a serialized database dump.

Three complete, offline, stdlib-first Python 3.11 implementations anchor the theory
in practice: a zero-framework REST parts-catalog service for the fictional
\emph{Northwind Robotics} universe that implements ETags, conditional requests,
content negotiation, and correct method semantics with nothing but
\texttt{http.server}; a URI design linter that flags style violations in route
tables; and a HAL-style hypermedia serializer. A worked OpenAPI 3.1 excerpt
\cite{openapi31} demonstrates the design-first workflow. The chapter closes with
thirty-six interview questions, five progressive capstone projects, and a
quick-reference card. Mastery objective: by the end of this chapter you should be
able to \emph{state}, \emph{defend}, and \emph{enforce} REST's constraints --- and
to recognize, with surgical precision, when an API is RESTful, when it is merely
HTTP, and why the difference matters for evolvability.

%% ----------------------------------------------------------------------------
\section{Learning Objectives \& Prerequisites}
\label{sec:ch3-objectives}

\subsection*{Learning Objectives}

After completing this chapter, its labs, and at least two capstone projects, you
will be able to:

\begin{enumerate}
  \item \textbf{Formalize} REST as an architectural style: state Fielding's six
        constraints verbatim-adjacent, name the properties each induces, and name
        what each forbids \cite{fielding2000rest}.
  \item \textbf{Diagnose} any HTTP API on the Richardson Maturity Model and
        articulate precisely which constraints it satisfies or violates
        \cite{richardson2013restful}.
  \item \textbf{Model} a business domain as resources: distinguish collections,
        documents, stores, and controllers; choose nouns over verbs; and design a
        coherent URI grammar (hierarchy, casing, pluralization, query versus path).
  \item \textbf{Apply} the uniform interface with precision: safety and
        idempotency of every method, correct status codes for each CRUD
        transition, conditional requests, and PUT versus PATCH semantics
        \cite{rfc9110}.
  \item \textbf{Evaluate} hypermedia (HATEOAS) as an engineering trade-off:
        compare HAL, JSON:API, and Collection+JSON, and decide when hypermedia
        pays for itself \cite{rfc8288}.
  \item \textbf{Design} representations for evolution: embedding versus linking,
        partial representations, and anti-corruption layering between the API and
        the database schema \cite{kleppmann2017ddia}.
  \item \textbf{Implement} constraint-compliant REST services with the Python
        standard library alone, and \textbf{govern} API style with linting,
        design-first OpenAPI workflows, and review gates \cite{openapi31,gehl2021apidesign}.
\end{enumerate}

\subsection*{Prerequisites}

This chapter assumes you have completed Chapters 0--2 of this book:

\begin{itemize}
  \item \textbf{Chapter 0} --- the course contract, environment setup, and the
        Northwind Robotics fictional universe used for all synthetic data.
  \item \textbf{Chapter 1} --- the HTTP wire protocol: requests, responses,
        headers, methods, status codes, and message bodies, as specified in
        RFC 9110 \cite{rfc9110}.
  \item \textbf{Chapter 2} --- consuming HTTP APIs from Python: \texttt{urllib}
        and \texttt{http.client}, parsing JSON, handling errors, and reading
        response headers.
\end{itemize}

You should be comfortable reading an HTTP request--response exchange in raw text
form and writing typed Python 3.11. No third-party packages are required for the
core listings; the optional route-table validator uses \texttt{pyyaml} when YAML
input is desired (Section~\ref{sec:ch3-dependencies}). Everything in this chapter
runs offline on a fresh Windows machine with PowerShell.

%% ----------------------------------------------------------------------------
\section{Deep Technical Mechanics \& Architectural Concepts}
\label{sec:ch3-mechanics}

\subsection{REST Is an Architectural Style, Not a Protocol}

The single most consequential sentence in Fielding's dissertation is also the most
ignored: REST is ``an architectural style for distributed hypermedia systems''
\cite{fielding2000rest}. An \emph{architectural style}, in Fielding's formal usage,
is a named, coordinated set of architectural constraints. A \emph{constraint} is a
rule that restricts the design space --- it \emph{removes} options. Styles earn
their keep because constraints, applied in combination, \emph{induce properties}:
emergent, measurable qualities of the resulting system such as scalability,
reliability, visibility, and evolvability.

\begin{definition}[Architectural Style]
An \emph{architectural style} is a coordinated set of constraints on the elements
(components, connectors, data) of an architecture and on the relationships among
them, chosen so that the constraints jointly induce a desired set of system
properties. REST is one such style; pipe-and-filter, event-driven integration, and
the layered monolith are others.
\end{definition}

Three practical consequences follow immediately:

\begin{enumerate}
  \item \textbf{REST is not HTTP.} HTTP is a transfer protocol whose design
        \emph{embodies} many REST constraints (Fielding co-authored it
        \cite{rfc9110}), which makes HTTP an unusually convenient substrate. But
        REST could in principle be realized over another transfer mechanism, and
        --- more importantly --- an HTTP API can violate every REST constraint
        while remaining perfectly valid HTTP.
  \item \textbf{``RESTful'' is gradated, not binary.} An API satisfies individual
        constraints to a degree. The honest engineering question is never ``is it
        REST?'' but ``which constraints does it honor, which does it trade away,
        and was that trade deliberate?''
  \item \textbf{Compliance is bought, not declared.} Each constraint has a cost
        (for example, statelessness moves session data onto the wire; cacheability
        risks staleness). The style is a negotiated exchange: constraints in,
        properties out.
\end{enumerate}

\begin{keyconceptbox}
REST is a \emph{style} (constraints $\rightarrow$ properties), not a protocol,
product, or specification. ``JSON over HTTP'' describes two data formats in
transit; it says nothing about architecture. An API is RESTful only insofar as it
honors Fielding's constraints --- and the most commonly skipped one is the uniform
interface's hypermedia clause.
\end{keyconceptbox}

\subsection{The Six Constraints, Formalized}
\label{sec:ch3-constraints}

Fielding derives REST incrementally, starting from the ``null style'' (no
constraints) and adding one constraint at a time \cite{fielding2000rest}. We follow
the same derivation, stating each constraint, its induced properties, and what it
forbids.

\subsubsection{Constraint 1: Client--Server}

\textbf{Statement.} Separate the user-interface concerns (client) from the
data-storage concerns (server) so that the two may evolve independently,
communicating only through a well-defined interface.

\textbf{Induced properties.} \emph{Portability} of the user interface across
platforms; \emph{scalability} by simplification of server components; and, most
valuably, \emph{independent evolvability} --- the server can change its storage
technology, and the client its rendering technology, without coordination.

\textbf{Forbids.} Shared-nothing violations in the small: the client may not reach
into server storage (no database links exposed to consumers), and the server may
not assume anything about client presentation.

\subsubsection{Constraint 2: Statelessness}

\begin{definition}[Statelessness (Fielding's sense)]
Each request from client to server must contain all of the information necessary
to understand the request, and cannot take advantage of any stored context on the
server. \emph{Session state is therefore kept entirely on the client.}
\end{definition}

Note the precise meaning: the server may of course hold \emph{resource} state (the
parts catalog itself). What it may not hold is \emph{session} state --- memory of
where a particular client is in an interaction. A request authenticated with a
bearer token, carrying its pagination cursor and all parameters, is stateless even
though the token validates against server-side data.

\textbf{Induced properties.} \emph{Visibility} (a monitoring system can understand
a request in isolation), \emph{reliability} (recovery from partial failure: any
request can simply be retried), and \emph{scalability} (no session affinity; any
replica can serve any request, and memory is freed between requests).

\textbf{Forbids.} Server-side session contexts, ``sticky'' conversational state,
and protocols whose requests are only interpretable relative to previous requests
on the same connection. The cost: repetitive data per request and reduced server
control over consistent application behavior --- the classic trade Fielding
acknowledges explicitly \cite{fielding2000rest}.

\subsubsection{Constraint 3: Cacheability}

\textbf{Statement.} The data within a response to a request must be implicitly or
explicitly labeled as cacheable or non-cacheable.

\textbf{Induced properties.} When a response is cacheable, a client cache may reuse
it for later, equivalent requests: \emph{efficiency}, \emph{scalability}, and
\emph{user-perceived performance} improve, sometimes by eliminating network
interactions entirely. HTTP realizes this constraint with validators
(\texttt{ETag}, \texttt{Last-Modified}) and freshness directives
(\texttt{Cache-Control}) \cite{rfc9110}.

\textbf{Forbids.} Ambiguity: a response that says nothing about its cacheability
forces intermediaries to guess, and guessed caching is how stale inventory gets
sold twice. The admitted risk is decreased \emph{reliability} if stale cached data
differs significantly from what a direct request would have returned.

\subsubsection{Constraint 4: Uniform Interface}

This is the constraint that distinguishes REST from generic HTTP usage, and it
decomposes into four sub-constraints \cite{fielding2000rest}:

\begin{enumerate}
  \item \textbf{Identification of resources.} Each resource is named by a stable
        identifier (in HTTP, a URI). The resource is the conceptual entity; the
        identifier is its name.
  \item \textbf{Manipulation of resources through representations.} Clients never
        touch the resource itself; they exchange \emph{representations} (e.g., a
        JSON document) of its current or desired state. A representation plus its
        metadata is sufficient to modify or delete the resource, given permission.
  \item \textbf{Self-descriptive messages.} Each message carries enough metadata
        (media type, method, status code, cache directives) to describe its own
        processing. Intermediaries can act on messages without private knowledge.
  \item \textbf{Hypermedia as the engine of application state (HATEOAS).} The
        client navigates the application by following links and forms discovered
        in representations --- not by constructing URIs from out-of-band
        knowledge.
\end{enumerate}

\begin{definition}[Resource]
A \emph{resource} is any named conceptual target of a hypertext reference: a
document, a temporal service (``today's weather in Lodz''), a collection of other
resources, a non-virtual object (a person), and so forth. A resource is a mapping
to a set of entities; the set may vary over time, provided the \emph{semantics}
of the mapping remain stable. ``The current price of part NX-1001'' is a fine
resource even though its value changes hourly.
\end{definition}

\begin{definition}[Representation]
A \emph{representation} is a sequence of bytes plus metadata (media type,
validators, links) that captures the current or intended state of a resource at
the time of a request. One resource admits many representations (JSON, CSV,
HTML); representations are negotiated, not assumed.
\end{definition}

\textbf{Induced properties.} \emph{Simplicity} of the overall architecture,
\emph{visibility} for intermediaries, and \emph{independent evolvability}, because
implementations are decoupled from the services they access. \textbf{The cost is
efficiency:} a uniform interface transfers information in a standardized form
rather than a form specialized to one application's needs.

\subsubsection{Constraint 5: Layered System}

\textbf{Statement.} The architecture is composed of hierarchical layers; each
component cannot ``see'' beyond the immediate layer with which it interacts.
Intermediaries --- gateways, load balancers, shared caches, content delivery
networks --- may be introduced at any layer without changing component semantics.

\textbf{Induced properties.} \emph{Scalability} (load balancing, shared caching at
organizational boundaries), \emph{security policy enforcement} at the edge, and
\emph{legacy encapsulation}. \textbf{The cost} is added latency and data overhead
per hop, which the cache constraint is expected to compensate for.

\subsubsection{Constraint 6: Code-on-Demand (Optional)}

\textbf{Statement.} The server may temporarily extend client functionality by
transferring executable code (scripts, applets) which the client executes.

\textbf{Induced properties.} \emph{Extensibility} and simplified clients (features
deploy without client installation). \textbf{Forbids/reduces:} \emph{visibility}
--- an intermediary cannot understand what arbitrary code will do --- which is
precisely why Fielding marks this constraint optional. JavaScript delivered to
browsers is its canonical realization; for machine-to-machine APIs it is almost
always omitted.

Figure~\ref{fig:ch3-layered} traces one conditional \texttt{GET} through a
concrete layered deployment, annotating which constraint each layer exploits.

\begin{figure}[h]
\centering
\begin{tikzpicture}[
  layer/.style={draw, rounded corners, minimum width=2.55cm, minimum height=1.05cm,
                align=center, font=\small, fill=blue!6},
  ann/.style={font=\scriptsize\itshape, text=packtgray, align=center},
  req/.style={->, thick, darkblue},
  res/.style={->, thick, dashed, packtgray}
]
  \node[layer] (client)  at (0,0)    {Client};
  \node[layer, right=1.15cm of client]  (edge) {Edge\\ gateway};
  \node[layer, right=1.15cm of edge]    (cache) {Shared\\ cache};
  \node[layer, right=1.15cm of cache]   (app)  {Catalog\\ service};
  \node[layer, right=1.15cm of app]     (store){Resource\\ store};

  \draw[req] (client) -- node[above, font=\scriptsize]
    {GET \texttt{/parts/NX-1001}} (edge);
  \draw[req] (edge) -- (cache);
  \draw[req] (cache) -- node[above, font=\scriptsize]
    {miss / revalidate} (app);
  \draw[req] (app) -- node[above, font=\scriptsize] {query} (store);

  \draw[res] (store.south) to[out=-35,in=-145]
    node[below, font=\scriptsize] {representation + \texttt{ETag}} (app.south);
  \draw[res] (cache.north) to[out=145,in=35]
    node[above, font=\scriptsize] {304 / cached hit} (edge.north);

  \node[ann, below=0.12cm of client] {session state\\ lives here};
  \node[ann, below=0.12cm of edge]   {layered system:\\ auth, TLS, routing};
  \node[ann, below=0.12cm of cache]  {cache constraint:\\ \texttt{ETag}, \texttt{Cache-Control}};
  \node[ann, below=0.12cm of app]    {stateless:\\ any replica\\ serves any request};
  \node[ann, below=0.12cm of store]  {resource state\\ (not session state)};
\end{tikzpicture}
\caption{Request lifecycle through a layered, constraint-compliant deployment.
Each layer sees only its immediate neighbors (layered system); the shared cache
short-circuits revalidation (cacheability); the service holds no session state
(statelessness).}
\label{fig:ch3-layered}
\end{figure}

Table~\ref{tab:ch3-constraints} summarizes the exchange each constraint offers.

\begin{table}[h]
  \centering
  \small
  \begin{tabularx}{\textwidth}{l X X}
    \toprule
    \textbf{Constraint} & \textbf{Induced properties} & \textbf{Cost / forbids} \\
    \midrule
    Client--server & portability, independent evolution & no shared storage access \\
    Statelessness & visibility, reliability, scalability & repeated context per request \\
    Cacheability & efficiency, perceived performance & staleness risk; must label responses \\
    Uniform interface & simplicity, visibility, decoupling & less efficient than bespoke forms \\
    Layered system & scalability, security at boundaries & per-hop latency \\
    Code-on-demand (opt.) & extensibility, thin clients & reduced visibility \\
    \bottomrule
  \end{tabularx}
  \caption{Fielding's six constraints as an exchange: constraints in, properties out \cite{fielding2000rest}.}
  \label{tab:ch3-constraints}
\end{table}

\subsection{``JSON over HTTP Is Not REST'': The Richardson Maturity Model}
\label{sec:ch3-rmm}

Because the industry adopted the word ``REST'' long before it absorbed Fielding's
constraints, a corrective vocabulary became necessary. Leonard Richardson's
\emph{Maturity Model} \cite{richardson2013restful} classifies HTTP APIs into four
levels according to which web technologies they actually exploit. It is best read
not as a moral hierarchy but as a diagnostic: each level names the constraints an
API has \emph{started} to honor. Figure~\ref{fig:ch3-rmm} shows the ladder.

\begin{figure}[h]
\centering
\begin{tikzpicture}[
  lvl/.style={draw, rounded corners, minimum width=7.4cm, minimum height=1.15cm,
              align=center, font=\small},
  desc/.style={font=\scriptsize\itshape, text=packtgray, align=left}
]
  \node[lvl, fill=red!12] (l0) at (0,0)
    {\textbf{Level 0 --- The Swamp of POX}\\ one URI, one method (POST), RPC payloads};
  \node[lvl, fill=orange!14, anchor=south west]
    at ($(l0.north east)+(1.0,0.5)$)
    {\textbf{Level 1 --- Resources}\\ many URIs (nouns), still one method};
  \node[lvl, fill=yellow!20, anchor=south west] (l2)
    at ($(l1.north east)+(1.0,0.5)$)
    {\textbf{Level 2 --- HTTP Verbs}\\ methods + status codes carry semantics};
  \node[lvl, fill=green!16, anchor=south west] (l3)
    at ($(l2.north east)+(1.0,0.5)$)
    {\textbf{Level 3 --- Hypermedia}\\ HATEOAS: representations carry controls};
  \draw[->, thick, packtgray] (l0.east) -- (l1.west);
  \draw[->, thick, packtgray] (l1.east) -- (l2.west);
  \draw[->, thick, packtgray] (l2.east) -- (l3.west);
  \node[desc, anchor=north west] at ($(l0.south west)+(0,-0.15)$)
    {RPC tunneling: HTTP as transport only};
  \node[desc, anchor=south east] at ($(l3.north east)+(0,0.15)$)
    {Full uniform interface \cite{fielding2000rest}};
\end{tikzpicture}
\caption{The Richardson Maturity Model. Each step adopts one more web
technology: resource URIs, then HTTP method/status semantics, then hypermedia
controls \cite{richardson2013restful}.}
\label{fig:ch3-rmm}
\end{figure}

\subsubsection{Level 0: The Swamp of POX}

A single endpoint receives every operation; the HTTP method and status code carry
no application meaning. This is HTTP used as a tunnel --- SOAP and many legacy
``XML-RPC over HTTP'' services live here.

\begin{minted}{text}
POST /rpc HTTP/1.1
Content-Type: application/json

{"action": "getPart", "sku": "NX-1001"}

HTTP/1.1 200 OK
Content-Type: application/json

{"status": "ok", "part": {"sku": "NX-1001", "name": "Servo Motor 12V"}}

POST /rpc HTTP/1.1
Content-Type: application/json

{"action": "getPart", "sku": "NOPE"}

HTTP/1.1 200 OK
Content-Type: application/json

{"status": "error", "message": "part not found"}
\end{minted}

Every failure is a \texttt{200 OK}; no proxy, cache, or monitoring tool can
distinguish success from failure, safe from unsafe. Nothing about the interaction
is visible to the web's infrastructure.

\subsubsection{Level 1: Resources}

The API introduces individually addressable resources --- URIs are nouns --- but
still tunnels all operations through one method.

\begin{minted}{text}
POST /parts/NX-1001 HTTP/1.1
Content-Type: application/json

{"action": "restock", "quantity": 40}

HTTP/1.1 200 OK
{"status": "ok", "stock": 180}
\end{minted}

Resources now exist as first-class addressable things (a genuine improvement:
URIs can be bookmarked, logged, and secured individually), but safety,
idempotency, and cacheability semantics remain unavailable because every request
is a \texttt{POST}.

\subsubsection{Level 2: HTTP Verbs}

Methods and status codes are used as RFC 9110 intends \cite{rfc9110}: \texttt{GET}
for retrieval (safe, cacheable), \texttt{PUT}/\texttt{PATCH} for update,
\texttt{DELETE} for removal, \texttt{POST} for subordinate creation, with status
codes that tell the truth.

\begin{minted}{text}
GET /parts/NX-1001 HTTP/1.1
Accept: application/json

HTTP/1.1 200 OK
ETag: "v3-9f2c"
Content-Type: application/json

{"sku": "NX-1001", "name": "Servo Motor 12V", "stock": 140}

GET /parts/NOPE HTTP/1.1

HTTP/1.1 404 Not Found
Content-Type: application/problem+json

{"type": "about:blank", "title": "Not Found", "status": 404,
 "detail": "No part with sku 'NOPE' exists."}
\end{minted}

Level 2 is where caches, conditional requests, idempotent retries, and honest
monitoring become possible. The overwhelming majority of well-run production APIs
live here.

\subsubsection{Level 3: Hypermedia Controls (HATEOAS)}

Representations carry \emph{controls} --- links and forms --- that advertise the
client's legal next steps, so application state is driven by server-supplied
hypermedia rather than client-side URI construction \cite{rfc8288}.

\begin{minted}{json}
{
  "sku": "NX-1001",
  "name": "Servo Motor 12V",
  "stock": 140,
  "_links": {
    "self":      { "href": "/parts/NX-1001" },
    "collection":{ "href": "/parts" },
    "category":  { "href": "/categories/actuators" },
    "restock":   { "href": "/parts/NX-1001/stock-adjustments",
                   "method": "POST",
                   "title": "Record a stock adjustment" }
  }
}
\end{minted}

A Level 3 client that needs to restock a part does not \emph{know} the adjustment
URI; it discovers it. The server can therefore restructure its URI space ---
insert a gateway, shard by region, rename collections --- without breaking
clients, because clients treat URIs as opaque (Section~\ref{sec:ch3-uri-grammar}).
Section~\ref{sec:ch3-hateoas} examines why this highest level remains rare and
when it is worth the investment.

\begin{warningbox}
Richardson's levels are diagnostic, not aspirational dogma. Level 3 buys
evolvability at real cost in tooling and client complexity. The engineering sin is
not stopping at Level 2; it is calling a Level 0 service ``RESTful'' and thereby
importing expectations (caches, safe retries, intermediaries) the service cannot
meet.
\end{warningbox}

%% ---------------------------------------------------------------------------
\subsection{Resource Modeling Discipline}
\label{sec:ch3-resources}

REST's central design act is deciding \emph{what the things are}. Get the resource
model right and methods, status codes, caching, and security policies fall into
place; get it wrong and no amount of framework sophistication will save the API.

\subsubsection{Nouns, Not Verbs}

A resource URI names a \emph{thing} (a noun); the HTTP method supplies the
\emph{operation} (the verb). The discipline exists because the uniform interface
already provides a fixed, universally understood verb vocabulary --- embedding
verbs in URIs duplicates that vocabulary with an unbounded, vendor-specific one,
destroying visibility: \texttt{GET /getPart?id=7} tells an intermediary nothing
about safety, while \texttt{GET /parts/NX-1001} declares itself safe and cacheable
by construction.

When a domain concept genuinely \emph{is} an action (``restock,'' ``authorize,''
``translate''), the RESTful resolution is to \emph{reify the action's result} as a
resource: instead of \texttt{POST /parts/NX-1001/restock}, create
\texttt{POST /parts/NX-1001/stock-adjustments} --- each adjustment is a record
with its own identity, audit trail, and lifecycle. The verb becomes a noun, and
the audit log you would have built anyway becomes the resource.

\subsubsection{The Four Resource Archetypes}

Following the taxonomy popularized by Mass\'e's \emph{REST API Design Rulebook}
(see Section~\ref{sec:ch3-bibliography}), resources come in four archetypes:

\begin{itemize}
  \item \textbf{Document} --- a singular, object-like resource:
        \texttt{/parts/NX-1001}. The REST analogue of a row or object instance.
  \item \textbf{Collection} --- a server-managed directory of documents:
        \texttt{/parts}. Clients \texttt{GET} it; \texttt{POST} proposes new
        members.
  \item \textbf{Store} --- a \emph{client}-managed resource repository: the
        client decides when an entry is created, updated, or deleted
        (\texttt{PUT /parts/NX-1001} with a client-chosen SKU is store-like).
  \item \textbf{Controller} --- the sanctioned escape hatch: a verb-like resource
        standing in for an operation that cannot be reified:
        \texttt{POST /catalog/exports} (creates an export job) is better modeled
        as a collection of \emph{jobs}; \texttt{POST /password-resets} similarly
        turns a procedure into a resource. Use controllers only when reification
        fails.
\end{itemize}

\subsubsection{Sub-Resources and Relationship Modeling}

Hierarchy in the path should express \emph{containment} or \emph{ownership} that
outlives either participant's individual addressing: a stock level exists only in
the context of a warehouse \emph{and} a part, so
\texttt{/warehouses/WH-01/stock/NX-1001} is defensible. Independent aggregates
with their own lifecycles deserve top-level URIs linked by hypermedia or foreign
identifiers, not deep nesting:

\begin{minted}{text}
GET /warehouses/WH-01                       # document
GET /warehouses/WH-01/stock                 # collection scoped by parent
GET /warehouses/WH-01/stock/NX-1001         # sub-resource (contained)
GET /purchase-orders?warehouse=WH-01        # relationship via query filter
GET /purchase-orders/PO-2026-0417           # independent aggregate, top level
\end{minted}

Rule of thumb, enforced by the linter in Listing~\ref{lst:ch3-urilint}: nest at
most two levels. Beyond that, prefer a top-level collection with a query filter.
Deep paths encode a \emph{single} navigational path into a graph that has many;
they ossify one query pattern and complicate authorization, caching keys, and
refactoring. Figure~\ref{fig:ch3-resources} shows the resource graph for the
Northwind Robotics warehouse domain developed in the lab
(Section~\ref{sec:ch3-lab}): solid edges are containment (path hierarchy),
dashed edges are associations (query filters or hypermedia links).

\begin{figure}[h]
\centering
\begin{tikzpicture}[
  coll/.style={draw, rounded corners, fill=green!12, minimum width=3.1cm,
               minimum height=0.85cm, font=\small\ttfamily, align=center},
  doc/.style={draw, rounded corners, fill=blue!8, minimum width=3.1cm,
              minimum height=0.85cm, font=\small\ttfamily, align=center},
  sub/.style={draw, rounded corners, fill=yellow!18, minimum width=3.6cm,
              minimum height=0.85cm, font=\small\ttfamily, align=center},
  owns/.style={->, thick, packtgray},
  assoc/.style={->, thick, dashed, darkblue}
]
  \node[coll] (warehouses)  at (0,0)     {/warehouses};
  \node[doc,   right=2.9cm of warehouses] (warehouse) {/warehouses/\{wid\}};
  \node[sub,   below=1.35cm of warehouse] (stock)
    {/warehouses/\{wid\}/stock};
  \node[coll,  below=1.35cm of warehouses] (parts) {/parts};
  \node[doc,   right=2.9cm of parts]      (part) {/parts/\{sku\}};
  \node[coll,  below=1.35cm of parts]     (pos) {/purchase-orders};
  \node[doc,   right=2.9cm of pos]        (po) {/purchase-orders/\{poid\}};

  \draw[owns] (warehouses) -- node[above, font=\scriptsize\itshape]{contains} (warehouse);
  \draw[owns] (warehouse) -- node[right, font=\scriptsize\itshape]{scopes} (stock);
  \draw[owns] (parts) -- node[above, font=\scriptsize\itshape]{contains} (part);
  \draw[owns] (pos) -- node[above, font=\scriptsize\itshape]{contains} (po);
  \draw[assoc] (stock.west) to[out=180,in=90]
    node[left, font=\scriptsize\itshape, align=center]{identifies part\\ (foreign key)} (parts.north);
  \draw[assoc] (po.north) to[out=90,in=-90]
    node[right, font=\scriptsize\itshape, align=center]{ships to\\ \texttt{?warehouse=}} (warehouse.south);
  \draw[assoc] (po.south) to[out=-90,in=0]
    node[below, font=\scriptsize\itshape, align=center]{orders parts\\ \texttt{?part=}} (part.south);
\end{tikzpicture}
\caption{Resource-relationship graph for the Northwind Robotics warehouse
domain. Containment (solid) appears in the path; independent aggregates
(purchase orders) get top-level URIs and relate via query filters and links
(dashed), keeping nesting at or below two levels.}
\label{fig:ch3-resources}
\end{figure}

\subsubsection{URI Design Grammar}
\label{sec:ch3-uri-grammar}

A URI decomposes as \texttt{scheme://authority/path?query\#fragment}; path
segments left-to-right descend a hierarchy. Production conventions, settled with
arguments rather than taste:

\begin{itemize}
  \item \textbf{Hierarchy.} Forward slashes express containment;
        \texttt{/warehouses/WH-01/stock} reads as a path \emph{into} the domain.
        Do not use hierarchy to encode \emph{queries}: filters belong after
        \texttt{?}.
  \item \textbf{Casing.} Lowercase \emph{kebab-case} for multi-word segments
        (\texttt{/purchase-orders}). URIs are case-sensitive per RFC 9110
        \cite{rfc9110}, but hosts and humans are not; camelCase paths
        (\texttt{/purchaseOrders}) produce duplicate-resource bugs behind
        case-insensitive proxies and normalize differently across gateways.
        Kebab-case is also what Zalando's and Microsoft's published guidelines
        converge on (Section~\ref{sec:ch3-bibliography}).
  \item \textbf{Pluralization.} The honest answer: the debate is noise; the
        invariant is \emph{consistency}. Plural collections win on three
        arguments: (1) the collection URI \texttt{/parts} and the document URI
        \texttt{/parts/NX-1001} then share a stem, making containment legible;
        (2) plural reads correctly for both \texttt{GET /parts} (many) and
        \texttt{GET /parts/NX-1001} (one member \emph{of} many); (3) singular
        resources remain expressible as singleton documents
        (\texttt{/catalog/profile}) without forcing every collection into
        grammatical contortions. Mixed schemes (\texttt{/part/NX-1001} but
        \texttt{/warehouses}) are the only unambiguously wrong answer, and the
        linter in Section~\ref{sec:ch3-code} treats them as defects.
  \item \textbf{Trailing slashes.} \texttt{/parts} and \texttt{/parts/} are
        \emph{different URIs} \cite{rfc9110}; permitting both doubles cache keys
        and splits analytics. Pick the bare form, and answer the slashed form
        with \texttt{301 Moved Permanently} --- a redirect is cacheable and
        self-healing; silent aliasing is not.
  \item \textbf{Query versus path.} Path segments identify; query parameters
        \emph{modulate}: filtering (\texttt{?category=sensors}), sorting
        (\texttt{?sort=-unitPrice}), pagination (\texttt{?cursor=...}), partial
        representation (\texttt{?fields=sku,name}). If removing the parameter
        changes \emph{which} resource is addressed, it belongs in the path; if it
        changes only \emph{how much} or \emph{which view} of that resource is
        returned, it belongs in the query.
\end{itemize}

\subsubsection{URI Opacity}

\begin{definition}[URI Opacity]
Clients must treat URIs as opaque identifiers: they may store, compare, and
dereference them, but must not parse, construct, or infer semantics from their
internal structure. Knowledge of how to \emph{form} a URI lives on the server and
is communicated through hypermedia or (pragmatically) a versioned contract such
as OpenAPI \cite{openapi31}.
\end{definition}

Opacity is what makes Level 3 restructuring safe and what makes even Level 2 APIs
evolvable: the day a client string-templates \texttt{/parts/} + sku, your URI
grammar has become an unchangeable public contract whether you intended it or
not. Fielding's formulation is blunt --- the only types of URI a RESTful system
may mint are those embedded in representations or constructed through
server-supplied templates \cite{fielding2000rest}.

%% ---------------------------------------------------------------------------
\subsection{Uniform Interface Semantics}
\label{sec:ch3-uniform}

\subsubsection{Safety and Idempotency}

RFC 9110 assigns every standard method two formal properties
\cite{rfc9110}. \emph{Safety} means the request is essentially read-only: the
client does not request, and cannot be held accountable for, any state change.
\emph{Idempotency} means the intended effect of multiple identical requests is
the same as that of a single request --- the property that makes automatic
retries of timed-out requests a sound engineering decision rather than a gamble.

\begin{table}[h]
  \centering
  \small
  \begin{tabularx}{\textwidth}{l c c X}
    \toprule
    \textbf{Method} & \textbf{Safe} & \textbf{Idempotent} & \textbf{Typical semantics on \texttt{/parts}} \\
    \midrule
    GET     & yes & yes & retrieve a representation; cacheable by default \\
    HEAD    & yes & yes & GET without the body; metadata probe \\
    OPTIONS & yes & yes & discover allowed methods (\texttt{Allow}) \\
    POST    & no  & no  & create subordinate; server assigns identity \\
    PUT     & no  & yes & full replacement (or creation) at a known URI \\
    PATCH   & no  & no\,\textsuperscript{*} & partial modification; semantics defined by media type \\
    DELETE  & no  & yes & remove; repeating yields 404, state unchanged \\
    \bottomrule
  \end{tabularx}
  \caption{Method properties per RFC 9110 \cite{rfc9110}.
  \textsuperscript{*}\,PATCH is not \emph{guaranteed} idempotent, but a specific
  patch document applied twice usually is --- design for it when you can.}
  \label{tab:ch3-methods}
\end{table}

Idempotency is defined over \emph{server state}, not responses: repeating
\texttt{DELETE /parts/NX-1001} returns 404 the second time, yet the method is
idempotent because the resulting state (no such part) is identical. This is why
load balancers and client libraries may transparently retry idempotent methods
and must never transparently retry \texttt{POST}.

\subsubsection{Status Codes for the CRUD Lifecycle}

Status codes are the server's half of the uniform interface: they let generic
components --- caches, gateways, monitors, retry loops --- react correctly
without private knowledge. Table~\ref{tab:ch3-status} gives the lifecycle map
used throughout this book \cite{rfc9110}.

\begin{table}[h]
  \centering
  \small
  \begin{tabularx}{\textwidth}{l l X}
    \toprule
    \textbf{Transition} & \textbf{Success} & \textbf{Failure semantics} \\
    \midrule
    Create (POST/PUT) & 201 + \texttt{Location} &
      400 invalid body; 409 conflict (duplicate SKU); 415 wrong media type; 422 semantically invalid \\
    Read (GET)        & 200 (304 on conditional hit) &
      404 unknown resource; 410 permanently retired; 406 no acceptable representation \\
    Replace (PUT)     & 200 or 204 (201 if created) &
      412 precondition failed (\texttt{If-Match} mismatch); 409 concurrent-edit conflict; 400/422 invalid \\
    Modify (PATCH)    & 200 or 204 &
      400 malformed patch document; 409 conflicting state; 422 patch not applicable; 415 wrong patch media type \\
    Remove (DELETE)   & 204 (200 with body) &
      404 already absent; 409 referenced by other resources \\
    Any               & --- &
      405 method not allowed (+ \texttt{Allow} header); 429 rate limited; 500/503 server-side \\
    \bottomrule
  \end{tabularx}
  \caption{Precise status codes for CRUD lifecycles. 404 in response to DELETE
  is informative, not an error the client must handle differently; 204 versus
  200 signals whether the response carries a representation.}
  \label{tab:ch3-status}
\end{table}

Two habits deserve special emphasis. First, \texttt{201 Created} without a
\texttt{Location} header is half an answer: the client learns \emph{that} a
resource was created but not \emph{where} it lives. Second, conditional requests
--- \texttt{If-None-Match} on GET (yielding \texttt{304 Not Modified}) and
\texttt{If-Match} on PUT/DELETE (yielding \texttt{412 Precondition Failed} on
mismatch) --- turn caching and optimistic concurrency from application-level
schemes into generic HTTP machinery that every intermediary already understands
\cite{rfc9110}. Listing~\ref{lst:ch3-catalog} implements both.

\subsubsection{PUT versus PATCH: Replacement versus Modification}

\texttt{PUT} is \emph{full replacement}: the enclosed representation becomes the
new state of the resource; omitted fields are reset, not preserved. It is
idempotent by construction. \texttt{PATCH} is \emph{partial modification}:
the body describes a \emph{change}, and its semantics are defined by the patch
media type. Two standardized patch formats dominate:

\begin{itemize}
  \item \textbf{JSON Merge Patch (RFC 7386).} The patch is itself a JSON
        document; it is recursively merged into the target: provided members
        replace, \texttt{null} members are \emph{removed}, absent members are
        untouched. Trivial to produce (any JSON object is a valid patch), but
        arrays are replaced wholesale and \texttt{null} can never be a stored
        value --- you cannot distinguish ``absent'' from ``set to null.''
\begin{minted}{json}
PATCH /parts/NX-1001        Content-Type: application/merge-patch+json
{ "unitPrice": { "currency": "USD", "amountCents": 2799 }, "legacyNote": null }
\end{minted}
  \item \textbf{JSON Patch (RFC 6902).} The patch is an \emph{array of
        operations} --- \texttt{add}, \texttt{remove}, \texttt{replace},
        \texttt{move}, \texttt{copy}, \texttt{test} --- applied in order with
        transaction semantics: if any operation fails, none is applied. The
        \texttt{test} operation doubles as a lightweight precondition
        (``apply only if \texttt{/stock} is still 140''). Verbose, but precise,
        array-aware, and atomically safe.
\begin{minted}{json}
PATCH /parts/NX-1001        Content-Type: application/json-patch+json
[
  { "op": "test",    "path": "/stock", "value": 140 },
  { "op": "replace", "path": "/stock", "value": 180 }
]
\end{minted}
\end{itemize}

Choose Merge Patch for forgiving, human-authored configuration updates; choose
JSON Patch for collaborative, concurrent editing where atomicity and
preconditions matter. Chapter 10 develops idempotency and concurrency control
in depth.

\subsubsection{Content Negotiation in Resource Design}

Because resources are manipulated \emph{through representations}, the same
resource may be rendered as JSON, CSV, or a vendor-specific media type. Server
driven negotiation uses the \texttt{Accept} header with quality values; the
server answers with the best available representation or \texttt{406 Not
Acceptable} \cite{rfc9110}:

\begin{minted}{text}
GET /parts HTTP/1.1
Accept: application/hal+json;q=1.0, text/csv;q=0.5

HTTP/1.1 200 OK
Content-Type: application/hal+json
Vary: Accept
\end{minted}

Design consequences: (1) always send \texttt{Vary: Accept} so shared caches key
variants separately; (2) negotiate \emph{format}, not \emph{semantics} --- a
\texttt{text/csv} rendering of \texttt{/parts} must describe the same resource
state, merely shaped differently; (3) media-type parameters
(\texttt{application/vnd.northwind.part+json;version=2}) are a legitimate
versioning surface, examined against URI versioning in Chapter 9.

%% ---------------------------------------------------------------------------
\subsection{HATEOAS and Hypermedia}
\label{sec:ch3-hateoas}

\subsubsection{Links and Affordances}

RFC 8288 standardizes \emph{web linking}: typed links between resources,
serializable in \texttt{Link} headers or in the body, each carrying a relation
type (\texttt{rel}) that names its semantics (\texttt{self}, \texttt{next},
\texttt{edit}, or an extension relation) \cite{rfc8288}. A link is a
\emph{navigation affordance}; a link plus method plus expected payload (HAL
forms, Collection+JSON templates) is an \emph{action affordance} --- the
hypermedia equivalent of an HTML form. Together they let the server, at runtime,
advertise the client's legal state transitions: the \texttt{restock} link in our
earlier example can simply be \emph{absent} when the part is discontinued, and
the client needs no `if discontinued` logic.

\subsubsection{Three Hypermedia Media Types Compared}

\begin{table}[h]
  \centering
  \small
  \begin{tabularx}{\textwidth}{l X X X}
    \toprule
    & \textbf{HAL} & \textbf{JSON:API} & \textbf{Collection+JSON} \\
    \midrule
    Media type & \texttt{application/hal+json} & \texttt{application/vnd.api+json} & \texttt{application/vnd.collection+json} \\
    Core model & \texttt{\_links}, \texttt{\_embedded}, curies & \texttt{data}/\texttt{attributes}/\texttt{relationships}, compound documents & collection of items + \texttt{queries} + write \texttt{template} \\
    Write support & out of scope (HAL Forms extension) & full CRUD semantics specified & template-driven POST/PUT bodies \\
    Philosophy & minimalism; bring your own conventions & opinionated full protocol (pagination, errors, sparse fields) & pure hypermedia; the representation \emph{is} the contract \\
    Best fit & adding links to an existing JSON API & greenfield APIs wanting everything standardized & reference designs, hypermedia research, generic clients \\
    \bottomrule
  \end{tabularx}
  \caption{Hypermedia media types for JSON APIs (see Section~\ref{sec:ch3-bibliography}).}
  \label{tab:ch3-hypermedia}
\end{table}

HAL is the pragmatic minimum: two reserved members buy link-driven navigation
with near-zero disruption \cite{rfc8288}. JSON:API standardizes nearly every
recurring design decision --- pagination, filtering, error objects, compound
documents --- which eliminates bikeshedding across teams but constrains resource
shapes. Collection+JSON is the purest realization of Fielding's constraint: its
\texttt{template} tells clients how to \emph{write} without any out-of-band
schema.

\subsubsection{Why Hypermedia Is Rare, and When It Pays Off}

Honesty about adoption matters for engineering judgment. Hypermedia APIs are
rare because (1) most API consumers are first-party web and mobile clients whose
release cadence matches the server's, so compile-time coupling to URIs is
tolerable; (2) mainstream tooling (OpenAPI generators, SDK pipelines) is built
around static contracts, not runtime-discovered affordances; (3) hypermedia
shifts complexity onto client authors, who must write generic link-followers
instead of calling endpoints. Hypermedia pays off when the cost asymmetry
flips: long-lived public APIs whose servers must restructure without breaking
integrations they do not control; workflow-heavy domains (order state machines,
approval chains) where the legal next actions genuinely vary per state; and
generic clients (crawlers, API browsers, agentic consumers) that cannot embed
per-API path knowledge. The professional default for this book: Level 2 with
\emph{selective} hypermedia --- \texttt{self}/\texttt{next}/\texttt{prev}
pagination links and state-transition links on workflow resources --- capturing
most of the evolvability at a fraction of the cost
\cite{richardson2013restful}.

%% ---------------------------------------------------------------------------
\subsection{Representation Design}
\label{sec:ch3-representations}

\subsubsection{Embedding versus Linking}

Should \texttt{GET /purchase-orders/PO-2026-0417} embed its line items, link to
them, or both? Embedding saves round trips and atomicity of read; linking saves
payload size, enables independent caching and freshness per resource, and keeps
the representation stable as the related collection grows. The disciplined
default: embed \emph{identity and summary} (id, name, and a link) of related
resources; offer full embedding behind an explicit opt-in
(\texttt{?embed=lines}) so the base representation stays small and cacheable.
Never embed unbounded collections in a document representation --- that is a
pagination accident waiting to happen (Chapter 8).

\subsubsection{Partial Representations}

Sparse fieldsets (\texttt{?fields=sku,name}) let bandwidth-sensitive clients
(mobile, high-fan-out aggregators) pay only for what they render. Keep the
mechanism simple, deterministic, and documented in the contract; projection must
never change \emph{which} resource is addressed --- that is what paths are for.

\subsubsection{API Is Not Database Schema: Anti-Corruption Thinking}

The fastest way to freeze your database forever is to serialize it. When
\texttt{GET /parts/NX-1001} is an \texttt{ORM} object passed through a JSON
encoder, every column rename, normalization, or storage migration becomes a
breaking API change --- and every breaking API change is now a schema change
negotiated with every consumer. The corrective is an \emph{anti-corruption
layer}: a deliberate translation between the storage model and the published
representation. The catalog in Section~\ref{sec:ch3-code} stores
\texttt{unit\_price\_cents} (an integer) but publishes a structured
\texttt{unitPrice} object with currency --- the representation is designed for
consumers, the storage column for accountants, and neither dictates the other.
Kleppmann's treatment of schema evolution gives the theoretical grounding:
systems evolve safely when writers and readers are decoupled by explicit
encodings with defined compatibility rules \cite{kleppmann2017ddia}.

\subsubsection{Evolution-Friendly Representations}

Four rules carry most of the weight \cite{gehl2021apidesign}: (1) \emph{additive
only} --- new optional fields are safe; removing or renaming fields is a
breaking change requiring a version transition (Chapter 9); (2) \emph{tolerant
readers} --- clients must ignore unknown members, and the contract must say so;
(3) \emph{never reuse names} --- a retired field's name stays retired, or old
clients will misread new data; (4) \emph{explicit nulls and enums} --- document
whether absent and null differ, and treat enum extension as a breaking change
for strict consumers. Applied consistently, these rules let a Level 2 API evolve
for years without version banners.

%% ---------------------------------------------------------------------------
\subsection{API as a Product}
\label{sec:ch3-product}

An API is a product whose customers are developers. Consumer-first design
inverts the usual implementation order: write the contract a consumer would
\emph{want} to integrate against \emph{before} writing the code that fulfills
it. The design-first workflow is concrete: draft the OpenAPI 3.1 document
\cite{openapi31}; review it with a representative consumer; lint it against the
organization's style guide; generate mocks so consumers can integrate against
the contract immediately; only then implement, with the contract driving tests
(Chapters 4 and 19). Consistency at scale is a governance problem, not a talent
problem: published style guides (Zalando, Microsoft --- see
Section~\ref{sec:ch3-bibliography}), automated linting in CI, and a review gate
for contract changes turn ``we value consistency'' from a sentiment into a
mechanism \cite{gehl2021apidesign}. The API Design Patterns corpus
\cite{gehl2021apidesign} catalogues the recurring decisions --- resource layout,
pagination tokens, long-running operations --- that such governance should
standardize once and reuse everywhere.

%% ---------------------------------------------------------------------------
\section{Production-Ready Code Implementation}
\label{sec:ch3-code}

The three listings in this section are complete, typed, and runnable offline on
Python 3.11+ with no third-party dependencies. All data is synthetic, drawn from
the fictional Northwind Robotics universe. Run each from PowerShell, e.g.
\texttt{python catalog\_api.py -{}-selftest}.

\subsection{A Stdlib-Only REST Mini-Framework: The Parts Catalog}

Listing~\ref{lst:ch3-catalog} implements the catalog domain on nothing but
\texttt{http.server}. It is deliberately framework-free so that every constraint
is visible in code: statelessness (no session state exists anywhere), the uniform
interface (method dispatch, status codes, \texttt{Allow} headers), cacheability
(ETags, conditional requests, \texttt{Cache-Control}), self-descriptive messages
(\texttt{Content-Type}, content negotiation, \texttt{Vary}), and layered-system
friendliness (it speaks plain HTTP intermediaries can cache and route).

\begin{minted}{python}
"""Northwind Robotics parts catalog API (stdlib-only, offline).

Demonstrates Fielding's constraints without frameworks:
  * uniform interface  - correct methods, status codes, Allow headers
  * cacheability       - ETag, If-None-Match (304), If-Match (412),
                         Cache-Control on every representation
  * statelessness      - zero session state; every request self-contained
  * self-descriptiveness - content negotiation (JSON / CSV), Vary: Accept

Python 3.11+. Run:
    python catalog_api.py --selftest     # deterministic in-process checks
    python catalog_api.py --port 8080    # serve; try curl.exe -i ...
"""
from __future__ import annotations

import argparse
import hashlib
import json
import logging
import re
from dataclasses import dataclass
from http import HTTPStatus
from http.server import BaseHTTPRequestHandler, ThreadingHTTPServer
from typing import Any
from urllib.parse import parse_qs, urlparse

logger = logging.getLogger("northwind.catalog")

# --------------------------------------------------------------------------
# Domain layer (anti-corruption: storage shape != representation shape)
# --------------------------------------------------------------------------

SKU_PATTERN = re.compile(r"^[A-Z]{2}-\d{4}$")
VALID_CATEGORIES = {"actuators", "sensors", "frames", "mobility"}


class ValidationError(Exception):
    """Raised when a representation fails boundary validation (400)."""


class ConflictError(Exception):
    """Raised when a write conflicts with existing state (409)."""


@dataclass
class Part:
    sku: str
    name: str
    category: str
    unit_price_cents: int
    stock: int
    version: int = 1

    def to_representation(self) -> dict[str, Any]:
        # Published shape is designed for consumers, not for the table:
        # money is a structured object; the storage column is an int.
        return {
            "sku": self.sku,
            "name": self.name,
            "category": self.category,
            "unitPrice": {
                "currency": "USD",
                "amountCents": self.unit_price_cents,
            },
            "stock": self.stock,
        }


def parse_part(body: bytes, *, sku_from_path: str | None = None) -> Part:
    """Validate a client representation into a Part. Fail fast, specifically."""
    try:
        data = json.loads(body.decode("utf-8"))
    except (UnicodeDecodeError, json.JSONDecodeError) as exc:
        raise ValidationError(f"request body is not valid JSON: {exc}") from exc
    if not isinstance(data, dict):
        raise ValidationError("request body must be a JSON object")

    price = data.get("unitPrice")
    if not isinstance(price, dict) or not isinstance(
        price.get("amountCents"), int
    ):
        raise ValidationError("unitPrice.amountCents (integer) is required")
    if price["amountCents"] < 0:
        raise ValidationError("unitPrice.amountCents must be >= 0")

    name = data.get("name")
    category = data.get("category")
    stock = data.get("stock")
    if not isinstance(name, str) or not name.strip():
        raise ValidationError("name (non-empty string) is required")
    if category not in VALID_CATEGORIES:
        raise ValidationError(
            f"category must be one of {sorted(VALID_CATEGORIES)}"
        )
    if not isinstance(stock, int) or stock < 0:
        raise ValidationError("stock (integer >= 0) is required")

    sku = sku_from_path or data.get("sku", "")
    if not SKU_PATTERN.match(sku):
        raise ValidationError("sku must match pattern AA-0000")
    if sku_from_path and data.get("sku") not in (None, sku_from_path):
        raise ValidationError("body sku must match the URI sku")

    return Part(
        sku=sku,
        name=name.strip(),
        category=category,
        unit_price_cents=price["amountCents"],
        stock=stock,
    )


class CatalogStore:
    """In-memory repository seeded with deterministic synthetic data."""

    def __init__(self) -> None:
        self._parts: dict[str, Part] = {}
        self._seed()

    def _seed(self) -> None:
        for part in (
            Part("NX-1001", "Servo Motor 12V", "actuators", 2599, 140),
            Part("NX-1002", "LIDAR Module R2", "sensors", 18999, 32),
            Part("NX-1003", "Aluminum Chassis Plate", "frames", 7450, 58),
            Part("NX-1004", "Omni Wheel 96mm", "mobility", 3250, 210),
        ):
            self._parts[part.sku] = part

    def list(self, category: str | None = None) -> list[Part]:
        parts = sorted(self._parts.values(), key=lambda p: p.sku)
        if category is not None:
            parts = [p for p in parts if p.category == category]
        return parts

    def get(self, sku: str) -> Part | None:
        return self._parts.get(sku)

    def create(self, part: Part) -> None:
        if part.sku in self._parts:
            raise ConflictError(f"part {part.sku} already exists")
        self._parts[part.sku] = part

    def replace(self, sku: str, part: Part) -> bool:
        """PUT semantics: full replacement. Returns True if it existed."""
        existed = sku in self._parts
        old = self._parts.get(sku)
        part.version = (old.version + 1) if old else 1
        self._parts[sku] = part
        return existed

    def delete(self, sku: str) -> bool:
        return self._parts.pop(sku, None) is not None


# --------------------------------------------------------------------------
# HTTP layer
# --------------------------------------------------------------------------

COLLECTION_RE = re.compile(r"^/parts$")
DOCUMENT_RE = re.compile(r"^/parts/(?P<sku>[A-Za-z0-9-]+)$")
JSON = "application/json"
CSV = "text/csv"


def etag_for(part: Part) -> str:
    digest = hashlib.sha256(
        json.dumps(part.to_representation(), sort_keys=True).encode()
    ).hexdigest()[:12]
    return f'"p{part.version}-{digest}"'


def negotiate(accept: str | None) -> str | None:
    """Pick the best supported media type, honoring q-values. None => 406."""
    if not accept or accept.strip() in ("", "*/*"):
        return JSON
    best: tuple[float, str] | None = None
    for item in accept.split(","):
        media, _, params = item.partition(";")
        media = media.strip().lower()
        quality = 1.0
        for param in params.split(";"):
            key, _, value = param.partition("=")
            if key.strip() == "q":
                try:
                    quality = float(value)
                except ValueError:
                    quality = 0.0
        if media in (JSON, CSV) and quality > 0:
            if best is None or quality > best[0]:
                best = (quality, media)
    return best[1] if best else None


def to_csv(reps: list[dict[str, Any]]) -> bytes:
    lines = ["sku,name,category,unit_price_cents,stock"]
    for r in reps:
        name = r["name"].replace('"', '""')
        lines.append(
            f'{r["sku"]},"{name}",{r["category"]},'
            f'{r["unitPrice"]["amountCents"]},{r["stock"]}'
        )
    return ("\r\n".join(lines) + "\r\n").encode("utf-8")


class CatalogHandler(BaseHTTPRequestHandler):
    server_version = "NorthwindCatalog/1.0"
    store: CatalogStore  # injected by make_server

    def log_message(self, fmt: str, *args: Any) -> None:
        logger.info("%s - %s", self.address_string(), fmt % args)

    # -- response helpers ---------------------------------------------------

    def _send(
        self,
        status: HTTPStatus,
        body: bytes = b"",
        content_type: str = JSON,
        extra: dict[str, str] | None = None,
    ) -> None:
        self.send_response(status)
        self.send_header("Content-Type", content_type)
        self.send_header("Content-Length", str(len(body)))
        self.send_header("Cache-Control", "no-cache")
        self.send_header("Vary", "Accept")
        for key, value in (extra or {}).items():
            self.send_header(key, value)
        self.end_headers()
        if body and self.command != "HEAD":
            self.wfile.write(body)

    def _problem(self, status: HTTPStatus, detail: str) -> None:
        body = json.dumps(
            {"title": status.phrase, "status": int(status), "detail": detail}
        ).encode()
        self._send(status, body, "application/problem+json")

    def _read_body(self) -> bytes:
        length = int(self.headers.get("Content-Length") or 0)
        return self.rfile.read(length) if length else b""

    def _check_json_content_type(self) -> None:
        ctype = (self.headers.get("Content-Type") or "").split(";")[0].strip()
        if ctype not in (JSON, "application/merge-patch+json"):
            raise UnsupportedMediaType(ctype)

    # -- method entry points ------------------------------------------------

    def do_GET(self) -> None:  # noqa: N802 (http.server naming)
        self._dispatch("GET")

    def do_HEAD(self) -> None:  # noqa: N802
        self._dispatch("GET")

    def do_POST(self) -> None:  # noqa: N802
        self._dispatch("POST")

    def do_PUT(self) -> None:  # noqa: N802
        self._dispatch("PUT")

    def do_PATCH(self) -> None:  # noqa: N802
        self._dispatch("PATCH")

    def do_DELETE(self) -> None:  # noqa: N802
        self._dispatch("DELETE")

    def do_OPTIONS(self) -> None:  # noqa: N802
        parsed = urlparse(self.path)
        allowed = (
            "GET, HEAD, POST, OPTIONS"
            if COLLECTION_RE.match(parsed.path)
            else "GET, HEAD, PUT, PATCH, DELETE, OPTIONS"
        )
        self._send(HTTPStatus.NO_CONTENT, extra={"Allow": allowed})

    # -- router ---------------------------------------------------------------

    def _dispatch(self, method: str) -> None:
        parsed = urlparse(self.path)
        path = parsed.path.rstrip("/") or "/"
        query = parse_qs(parsed.query)
        doc = DOCUMENT_RE.match(path)
        try:
            if COLLECTION_RE.match(path):
                self._collection(method, query)
            elif doc:
                self._document(method, doc.group("sku"))
            else:
                self._problem(HTTPStatus.NOT_FOUND, f"no resource at {path}")
        except ValidationError as exc:
            self._problem(HTTPStatus.BAD_REQUEST, str(exc))
        except ConflictError as exc:
            self._problem(HTTPStatus.CONFLICT, str(exc))
        except UnsupportedMediaType as exc:
            self._problem(
                HTTPStatus.UNSUPPORTED_MEDIA_TYPE,
                f"unsupported media type: {exc}",
            )
        except Exception:  # never leak internals over the wire
            logger.exception("unhandled error processing %s %s", method, path)
            self._problem(HTTPStatus.INTERNAL_SERVER_ERROR, "internal error")

    def _reject(self, allowed: str) -> None:
        body = json.dumps(
            {
                "title": "Method Not Allowed",
                "status": int(HTTPStatus.METHOD_NOT_ALLOWED),
                "detail": "method not allowed on this resource",
            }
        ).encode()
        self._send(
            HTTPStatus.METHOD_NOT_ALLOWED,
            body,
            "application/problem+json",
            extra={"Allow": allowed},
        )

    # -- /parts (collection) --------------------------------------------------

    def _collection(self, method: str, query: dict[str, list[str]]) -> None:
        if method == "GET":
            category = query.get("category", [None])[0]
            reps = [p.to_representation() for p in self.store.list(category)]
            media = negotiate(self.headers.get("Accept"))
            if media is None:
                self._problem(
                    HTTPStatus.NOT_ACCEPTABLE, "supported: application/json, text/csv"
                )
                return
            body = (
                json.dumps({"parts": reps, "count": len(reps)}).encode()
                if media == JSON
                else to_csv(reps)
            )
            self._send(HTTPStatus.OK, body, media)
        elif method == "POST":
            self._check_json_content_type()
            part = parse_part(self._read_body())
            self.store.create(part)
            body = json.dumps(part.to_representation()).encode()
            self._send(
                HTTPStatus.CREATED,
                body,
                extra={
                    "Location": f"/parts/{part.sku}",
                    "ETag": etag_for(part),
                },
            )
        else:
            self._reject("GET, HEAD, POST, OPTIONS")

    # -- /parts/{sku} (document) ------------------------------------------------

    def _document(self, method: str, sku: str) -> None:
        part = self.store.get(sku)
        if method == "GET":  # HEAD is dispatched here with body suppressed
            if part is None:
                self._problem(HTTPStatus.NOT_FOUND, f"no part with sku {sku!r}")
                return
            etag = etag_for(part)
            if self.headers.get("If-None-Match") == etag:
                self._send(HTTPStatus.NOT_MODIFIED, extra={"ETag": etag})
                return
            media = negotiate(self.headers.get("Accept"))
            if media is None:
                self._problem(HTTPStatus.NOT_ACCEPTABLE, "supported: json, csv")
                return
            rep = part.to_representation()
            body = (
                json.dumps(rep).encode() if media == JSON else to_csv([rep])
            )
            self._send(HTTPStatus.OK, body, media, extra={"ETag": etag})
        elif method == "PUT":
            self._check_json_content_type()
            if part and self.headers.get("If-Match") not in (None, etag_for(part)):
                self._problem(
                    HTTPStatus.PRECONDITION_FAILED,
                    "If-Match does not match current ETag; re-read and retry",
                )
                return
            new_part = parse_part(self._read_body(), sku_from_path=sku)
            existed = self.store.replace(sku, new_part)
            etag = etag_for(self.store.get(sku))  # type: ignore[arg-type]
            status = HTTPStatus.OK if existed else HTTPStatus.CREATED
            extra = {"ETag": etag}
            if not existed:
                extra["Location"] = f"/parts/{sku}"
            self._send(status, b"", extra=extra)
        elif method == "PATCH":
            if part is None:
                self._problem(HTTPStatus.NOT_FOUND, f"no part with sku {sku!r}")
                return
            ctype = (self.headers.get("Content-Type") or "").split(";")[0]
            if ctype.strip() != "application/merge-patch+json":
                raise UnsupportedMediaType(ctype or "<missing>")
            patch = json.loads(self._read_body().decode("utf-8") or "{}")
            merged = merge_patch(part.to_representation(), patch)
            updated = parse_part(
                json.dumps(merged).encode(), sku_from_path=sku
            )
            self.store.replace(sku, updated)  # bumps the version counter
            stored = self.store.get(sku)
            assert stored is not None  # just written by replace()
            body = json.dumps(stored.to_representation()).encode()
            self._send(HTTPStatus.OK, body, extra={"ETag": etag_for(stored)})
        elif method == "DELETE":
            if self.store.delete(sku):
                self._send(HTTPStatus.NO_CONTENT)
            else:
                self._problem(HTTPStatus.NOT_FOUND, f"no part with sku {sku!r}")
        else:
            self._reject("GET, HEAD, PUT, PATCH, DELETE, OPTIONS")


class UnsupportedMediaType(Exception):
    """Raised when Content-Type is not one the resource understands (415)."""


def merge_patch(target: dict[str, Any], patch: dict[str, Any]) -> dict[str, Any]:
    """RFC 7386 JSON Merge Patch: merge keys; null deletes the key."""
    result = dict(target)
    for key, value in patch.items():
        if value is None:
            result.pop(key, None)
        elif isinstance(value, dict) and isinstance(result.get(key), dict):
            result[key] = merge_patch(result[key], value)
        else:
            result[key] = value
    return result


def make_server(port: int) -> ThreadingHTTPServer:
    CatalogHandler.store = CatalogStore()
    return ThreadingHTTPServer(("127.0.0.1", port), CatalogHandler)


def run_selftest() -> None:
    """Deterministic in-process checks; no sockets, no network."""
    from http.client import HTTPConnection
    import threading

    server = make_server(0)
    thread = threading.Thread(target=server.serve_forever, daemon=True)
    thread.start()
    host, port = server.server_address
    conn = HTTPConnection(host, port, timeout=5)
    try:
        conn.request("GET", "/parts/NX-1001")
        resp = conn.getresponse()
        assert resp.status == 200 and resp.getheader("ETag")
        etag = resp.getheader("ETag")
        resp.read()

        conn.request("GET", "/parts/NX-1001", headers={"If-None-Match": etag})
        assert conn.getresponse().status == 304

        conn.request("GET", "/parts/NOPE")
        assert conn.getresponse().status == 404

        conn.request(
            "POST",
            "/parts",
            body=json.dumps({
                "sku": "NX-1005", "name": "Torque Wrench Arm",
                "category": "actuators",
                "unitPrice": {"currency": "USD", "amountCents": 9900},
                "stock": 12,
            }),
            headers={"Content-Type": "application/json"},
        )
        resp = conn.getresponse()
        assert resp.status == 201
        assert resp.getheader("Location") == "/parts/NX-1005"
        resp.read()

        conn.request("DELETE", "/parts/NX-1005")
        assert conn.getresponse().status == 204
        conn.request("DELETE", "/parts/NX-1005")
        assert conn.getresponse().status == 404

        conn.request(
            "GET", "/parts", headers={"Accept": "application/xml"}
        )
        assert conn.getresponse().status == 406
        print("selftest: all checks passed (200/304/404/201/204/404/406)")
    finally:
        conn.close()
        server.shutdown()
        server.server_close()


if __name__ == "__main__":
    parser = argparse.ArgumentParser(description=__doc__)
    parser.add_argument("--port", type=int, default=8080)
    parser.add_argument("--selftest", action="store_true")
    args = parser.parse_args()
    logging.basicConfig(level=logging.INFO, format="%(levelname)s %(message)s")
    if args.selftest:
        run_selftest()
    else:
        logger.info("serving Northwind parts catalog on 127.0.0.1:%d", args.port)
        make_server(args.port).serve_forever()
\end{minted}

\noindent\textbf{Listing 3.1 --- \texttt{catalog\_api.py}: constraint-compliant
REST service on \texttt{http.server}.}\label{lst:ch3-catalog}

Note what the framework-free implementation makes unavoidable: the \texttt{ETag}
is derived from representation bytes plus a version counter; \texttt{If-Match} on
\texttt{PUT} implements optimistic concurrency for free; \texttt{405} responses
carry \texttt{Allow}; \texttt{415} guards the content-type boundary; and the
\texttt{Cache-Control: no-cache} plus \texttt{ETag} combination means caches must
revalidate but may reuse bodies --- correctness without giving up the network
savings \cite{rfc9110}. The one deliberate simplification is storage: the
in-memory store stands in for a repository layer that Chapter 6 replaces with a
real persistence stack behind the identical interface.

\subsection{A URI Design Validator}

Style guides die in wikis unless they are executable. Listing~\ref{lst:ch3-urilint}
implements the URI grammar rules of Section~\ref{sec:ch3-uri-grammar} as a linter
over a route table --- a JSON (stdlib) or YAML (requires \texttt{pyyaml},
Section~\ref{sec:ch3-dependencies}) document shaped like a simplified OpenAPI
\texttt{paths} object. It exits non-zero when it finds violations, so it can gate
CI exactly like a code linter.

\begin{minted}{python}
"""uri_lint.py - enforce the URI design grammar on a route table.

Usage (PowerShell):
    python uri_lint.py                    # lint the embedded sample table
    python uri_lint.py routes.json        # JSON route table (stdlib only)
    python uri_lint.py routes.yaml        # YAML route table (needs pyyaml)

Exit code 0 = clean, 1 = violations found, 2 = input could not be parsed.
Deterministic: findings are sorted; no clock, no network, no randomness.
"""
from __future__ import annotations

import json
import logging
import re
import sys
from dataclasses import dataclass
from pathlib import Path
from typing import Any

logger = logging.getLogger("northwind.uri_lint")

MAX_NESTING_DEPTH = 2  # collection/document pairs, per style guide

# Verb lexicon: segments that signal RPC-style tunneling (rule R001).
VERB_LEXICON = {
    "get", "fetch", "read", "list", "query", "search", "find",
    "create", "add", "insert", "register", "make", "new",
    "update", "edit", "modify", "change", "set", "patch",
    "delete", "remove", "destroy", "purge",
    "do", "process", "execute", "run", "submit", "trigger",
    "calculate", "compute", "export", "import", "sync",
}

KEBAB_RE = re.compile(r"^[a-z0-9]+(-[a-z0-9]+)*$")
PARAM_RE = re.compile(r"^\{[a-zA-Z][a-zA-Z0-9_]*\}$")


@dataclass(frozen=True)
class Finding:
    rule: str
    path: str
    message: str

    def render(self) -> str:
        return f"[{self.rule}] {self.path}: {self.message}"


def static_segments(path: str) -> list[str]:
    """Path segments minus template parameters like {sku}."""
    return [s for s in path.split("/") if s and not PARAM_RE.match(s)]


def lint_path(path: str) -> list[Finding]:
    findings: list[Finding] = []

    if path != "/" and path.endswith("/"):                       # R005
        findings.append(Finding(
            "R005", path,
            "trailing slash creates a duplicate URI; use the bare form "
            "and 301-redirect the slashed form",
        ))

    segments = [s for s in path.split("/") if s]
    for seg in segments:
        if PARAM_RE.match(seg):
            continue
        if seg.lower() in VERB_LEXICON:                          # R001
            findings.append(Finding(
                "R001", path,
                f"verb {seg!r} in path; the HTTP method supplies the verb - "
                f"reify the action as a resource instead",
            ))
        if not KEBAB_RE.match(seg):                              # R004
            findings.append(Finding(
                "R004", path,
                f"segment {seg!r} is not lowercase kebab-case "
                f"(no camelCase, snake_case, or mixed case)",
            ))

    depth = (len(segments) + 1) // 2                             # R002
    if depth > MAX_NESTING_DEPTH:
        findings.append(Finding(
            "R002", path,
            f"nesting depth {depth} exceeds {MAX_NESTING_DEPTH}; "
            f"prefer a top-level collection with a query filter",
        ))
    return findings


def lint_pluralization(paths: list[str]) -> list[Finding]:
    """R003: collection segments must be consistently plural."""
    singular: list[str] = []
    plural: list[str] = []
    for path in paths:
        segs = static_segments(path)
        if not segs:
            continue
        first = segs[0]
        (plural if first.endswith("s") else singular).append(first)
    if singular and plural:
        return [Finding(
            "R003", ", ".join(sorted(set(singular))),
            "singular collection segment(s) mixed with plural "
            f"{sorted(set(plural))}; pick plural everywhere",
        )]
    return []


def lint_route_table(table: dict[str, Any]) -> list[Finding]:
    paths = table.get("paths")
    if not isinstance(paths, dict) or not paths:
        raise ValueError("route table must contain a non-empty 'paths' object")
    findings: list[Finding] = []
    for path in paths:
        if not isinstance(path, str) or not path.startswith("/"):
            raise ValueError(f"invalid path key: {path!r}")
        findings.extend(lint_path(path))
    findings.extend(lint_pluralization(list(paths)))
    return sorted(findings, key=lambda f: (f.path, f.rule))


def load_table(path: Path) -> dict[str, Any]:
    text = path.read_text(encoding="utf-8")
    if path.suffix in (".yaml", ".yml"):
        try:
            import yaml  # type: ignore[import-untyped]
        except ImportError as exc:
            raise RuntimeError(
                "YAML input requires pyyaml: pip install pyyaml"
            ) from exc
        data = yaml.safe_load(text)
    else:
        data = json.loads(text)
    if not isinstance(data, dict):
        raise ValueError("route table must be a JSON/YAML object")
    return data


SAMPLE_TABLE: dict[str, Any] = {
    "paths": {
        "/parts": {},
        "/parts/{sku}": {},
        "/parts/{sku}/restock": {},    # action-like; reify as stock-adjustments
        "/getPart": {},                # R003 + R004 (getPart: camelCase, singular)
        "/warehouses/{wid}/stock/{sku}/history": {},  # R002: depth 3
        "/purchaseOrders/{id}": {},            # R004: camelCase
        "/suppliers/": {},                     # R005: trailing slash
    }
}


def main(argv: list[str]) -> int:
    logging.basicConfig(level=logging.WARNING, format="%(levelname)s %(message)s")
    if len(argv) > 1:
        try:
            table = load_table(Path(argv[1]))
        except (OSError, ValueError, RuntimeError,
                json.JSONDecodeError) as exc:
            logger.error("cannot load route table: %s", exc)
            return 2
    else:
        table = SAMPLE_TABLE

    findings = lint_route_table(table)
    for finding in findings:
        print(finding.render())
    print(f"\n{len(findings)} violation(s) in {len(table['paths'])} path(s)")
    return 1 if findings else 0


if __name__ == "__main__":
    sys.exit(main(sys.argv))
\end{minted}

\noindent\textbf{Listing 3.2 --- \texttt{uri\_lint.py}: executable URI style
guide.}\label{lst:ch3-urilint}

Run against its embedded sample, the linter reports five findings across seven
paths (R002--R005 all fire) --- style defects caught mechanically, before a
single line of service code exists. The lexicon deliberately contains only
generic verbs; domain-specific action words like \texttt{restock} are caught by
human review guided by the reification rule, not by an ever-growing word list. This is design governance as code: the same
rule IDs can be referenced in review comments and waivers.

\subsection{A HAL-Style Hypermedia Serializer}

Listing~\ref{lst:ch3-hal} closes the loop on Section~\ref{sec:ch3-hateoas}: it
renders the catalog domain as HAL-like documents with \texttt{\_links} and
\texttt{\_embedded}, keeping serialization separate from both storage and
transport.

\begin{minted}{python}
"""hal_serializer.py - render catalog resources as HAL-like documents.

Stdlib-only and deterministic (sorted keys, fixed sample data). The point:
links are DATA emitted by the server, not strings concatenated by clients.
"""
from __future__ import annotations

import json
from dataclasses import dataclass, field
from typing import Any

BASE = "https://api.northwind-robotics.example"  # fictional host, never called


@dataclass(frozen=True)
class Link:
    href: str
    method: str = "GET"
    title: str | None = None
    templated: bool = False

    def render(self) -> dict[str, Any]:
        out: dict[str, Any] = {"href": self.href}
        if self.method != "GET":
            out["method"] = self.method
        if self.title:
            out["title"] = self.title
        if self.templated:
            out["templated"] = True
        return out


@dataclass
class HalResource:
    state: dict[str, Any]
    links: dict[str, Link | list[Link]] = field(default_factory=dict)
    embedded: dict[str, list["HalResource"]] = field(default_factory=dict)

    def to_dict(self) -> dict[str, Any]:
        doc: dict[str, Any] = dict(self.state)
        doc["_links"] = {
            rel: ([lnk.render() for lnk in link]
                  if isinstance(link, list) else link.render())
            for rel, link in sorted(self.links.items())
        }
        if self.embedded:
            doc["_embedded"] = {
                rel: [res.to_dict() for res in resources]
                for rel, resources in sorted(self.embedded.items())
            }
        return doc


def part_resource(part: dict[str, Any]) -> HalResource:
    sku = part["sku"]
    links: dict[str, Link | list[Link]] = {
        "self": Link(f"{BASE}/parts/{sku}"),
        "collection": Link(f"{BASE}/parts", title="All parts"),
        "category": Link(f"{BASE}/categories/{part['category']}"),
    }
    # Action affordances are state-dependent: only in-stock, active parts
    # advertise the restock transition. Clients need no business rule for this.
    if part["stock"] > 0 and part.get("status", "active") == "active":
        links["restock"] = Link(
            f"{BASE}/parts/{sku}/stock-adjustments",
            method="POST",
            title="Record a stock adjustment",
        )
    return HalResource(state={k: v for k, v in part.items() if k != "status"},
                       links=links)


def parts_collection(parts: list[dict[str, Any]], *, cursor: str | None,
                     next_cursor: str | None) -> HalResource:
    collection = HalResource(
        state={"count": len(parts)},
        links={"self": Link(f"{BASE}/parts" + (f"?cursor={cursor}" if cursor
                                               else ""))},
        embedded={"item": [part_resource(p) for p in parts]},
    )
    # Pagination links are the minimum viable hypermedia: the server owns the
    # cursor format, so it can change storage without breaking clients.
    if next_cursor is not None:
        collection.links["next"] = Link(f"{BASE}/parts?cursor={next_cursor}")
    return collection


CATALOG_SAMPLE: list[dict[str, Any]] = [
    {"sku": "NX-1001", "name": "Servo Motor 12V", "category": "actuators",
     "unitPrice": {"currency": "USD", "amountCents": 2599}, "stock": 140},
    {"sku": "NX-1002", "name": "LIDAR Module R2", "category": "sensors",
     "unitPrice": {"currency": "USD", "amountCents": 18999}, "stock": 32},
    {"sku": "NX-1003", "name": "Aluminum Chassis Plate", "category": "frames",
     "unitPrice": {"currency": "USD", "amountCents": 7450}, "stock": 0,
     "status": "discontinued"},
]


def main() -> None:
    document = part_resource(CATALOG_SAMPLE[0]).to_dict()
    listing = parts_collection(
        CATALOG_SAMPLE, cursor=None, next_cursor="c3RlcDoy"
    ).to_dict()
    print("--- document ---")
    print(json.dumps(document, indent=2, sort_keys=True))
    print("--- collection ---")
    print(json.dumps(listing, indent=2, sort_keys=True))


if __name__ == "__main__":
    main()
\end{minted}

\noindent\textbf{Listing 3.3 --- \texttt{hal\_serializer.py}: HAL-style
serializer with state-dependent action affordances.}\label{lst:ch3-hal}

Two details reward attention. First, the \texttt{restock} affordance is
\emph{computed from resource state}: a discontinued part simply stops
advertising the transition, and a hypermedia-driven client disables the feature
without a code change. Second, the pagination \texttt{next} link is an opaque
cursor the server owns --- the client stores and dereferences it but never
constructs it, which is URI opacity (Section~\ref{sec:ch3-uri-grammar}) doing
real work.

\subsection{Design-First Contract: OpenAPI 3.1 Excerpt}

Section~\ref{sec:ch3-product} argued for contract-first development.
Listing~\ref{lst:ch3-openapi} is the worked excerpt for the catalog; the prose
commentary follows in the block's design notes.

\begin{minted}{yaml}
openapi: 3.1.0
info:
  title: Northwind Robotics Parts Catalog API
  version: 1.0.0
  summary: Resource-oriented catalog of robot parts (synthetic training data).
  license: { name: Proprietary }
jsonSchemaDialect: https://json-schema.org/draft/2020-12/schema
servers:
  - url: http://localhost:8080
    description: Offline development server (Listing 3.1)
tags:
  - name: parts
    description: Catalog part documents and the parts collection.

paths:
  /parts:
    get:
      summary: List parts
      operationId: listParts
      tags: [parts]
      parameters:
        - name: category
          in: query
          description: Filter by category. Query modulates; path identifies.
          schema:
            type: string
            enum: [actuators, sensors, frames, mobility]
      responses:
        "200":
          description: A page of part representations.
          headers:
            Cache-Control:
              schema: { type: string }
          content:
            application/json:
              schema: { $ref: "#/components/schemas/PartList" }
            text/csv:
              schema: { type: string }
        "406": { $ref: "#/components/responses/NotAcceptable" }
    post:
      summary: Create a part (server assigns nothing; client supplies SKU)
      operationId: createPart
      tags: [parts]
      requestBody:
        required: true
        content:
          application/json:
            schema: { $ref: "#/components/schemas/PartInput" }
      responses:
        "201":
          description: Created; Location header addresses the new document.
          headers:
            Location: { schema: { type: string } }
            ETag: { schema: { type: string } }
        "400": { $ref: "#/components/responses/BadRequest" }
        "409": { $ref: "#/components/responses/Conflict" }
        "415": { $ref: "#/components/responses/UnsupportedMediaType" }

  /parts/{sku}:
    parameters:
      - name: sku
        in: path
        required: true
        schema: { type: string, pattern: "^[A-Z]{2}-[0-9]{4}$" }
    get:
      summary: Retrieve a part; supports conditional requests
      operationId: getPart
      tags: [parts]
      parameters:
        - name: If-None-Match
          in: header
          schema: { type: string }
      responses:
        "200":
          description: The part representation, with validators.
          headers:
            ETag: { schema: { type: string } }
          content:
            application/json:
              schema: { $ref: "#/components/schemas/Part" }
        "304": { description: Not modified; re-use the cached copy. }
        "404": { $ref: "#/components/responses/NotFound" }
    put:
      summary: Full replacement (idempotent); If-Match guards lost updates
      operationId: replacePart
      tags: [parts]
      parameters:
        - name: If-Match
          in: header
          schema: { type: string }
      requestBody:
        required: true
        content:
          application/json:
            schema: { $ref: "#/components/schemas/PartInput" }
      responses:
        "200": { description: Replaced. }
        "201": { description: Created at a client-known URI (store semantics). }
        "412": { description: ETag precondition failed; re-read and retry. }
    patch:
      summary: Partial update via JSON Merge Patch (RFC 7386)
      operationId: mergePatchPart
      tags: [parts]
      requestBody:
        required: true
        content:
          application/merge-patch+json:
            schema: { type: object }
      responses:
        "200": { description: Updated representation returned. }
        "415": { $ref: "#/components/responses/UnsupportedMediaType" }
    delete:
      summary: Remove a part (idempotent effect)
      operationId: deletePart
      tags: [parts]
      responses:
        "204": { description: Removed; no body. }
        "404": { $ref: "#/components/responses/NotFound" }

components:
  schemas:
    Money:
      type: object
      required: [currency, amountCents]
      properties:
        currency: { type: string, minLength: 3, maxLength: 3 }
        amountCents: { type: integer, minimum: 0 }
      additionalProperties: false
    PartInput:
      type: object
      required: [sku, name, category, unitPrice, stock]
      properties:
        sku: { type: string, pattern: "^[A-Z]{2}-[0-9]{4}$" }
        name: { type: string, minLength: 1 }
        category: { type: string, enum: [actuators, sensors, frames, mobility] }
        unitPrice: { $ref: "#/components/schemas/Money" }
        stock: { type: integer, minimum: 0 }
      additionalProperties: false
    Part:
      allOf:
        - $ref: "#/components/schemas/PartInput"
      description: Server representation; identical shape here by choice,
        but the contract (not the database) is the source of truth.
    PartList:
      type: object
      required: [parts, count]
      properties:
        parts:
          type: array
          items: { $ref: "#/components/schemas/Part" }
        count: { type: integer }
  responses:
    BadRequest:
      description: Malformed or invalid representation.
      content:
        application/problem+json:
          schema: { $ref: "#/components/schemas/Problem" }
    NotFound: { description: Unknown resource. }
    Conflict: { description: State conflict (duplicate SKU, stale edit). }
    NotAcceptable: { description: No supported representation matches Accept. }
    UnsupportedMediaType: { description: Content-Type not supported. }
  # Problem schema elided for brevity: title/status/detail per problem+json.
\end{minted}

\noindent\textbf{Listing 3.4 --- \texttt{openapi.yaml} (excerpt): design-first
contract for the catalog.}\label{lst:ch3-openapi}

Design commentary: every response declares its headers (\texttt{ETag},
\texttt{Location}) because headers are contract, not decoration; error responses
are reusable components so the \texttt{400}s across resources can never drift
apart; \texttt{additionalProperties: false} makes the contract closed by default,
forcing field additions through deliberate review (the evolution discipline of
Section~\ref{sec:ch3-representations}); and conditional-request headers are
first-class parameters, so generated SDKs expose optimistic concurrency to every
consumer. Chapter 4 develops serialization contracts; Chapter 9 extends this
document with a versioning strategy.

%% ---------------------------------------------------------------------------
\section{Common Pitfalls \& Anti-Patterns}
\label{sec:ch3-pitfalls}

Each anti-pattern below is stated with its symptom, the constraint it breaks,
and the corrective. All are detected by Listing~\ref{lst:ch3-urilint} where they
appear in URIs.

\begin{enumerate}
  \item \textbf{Verbs in URIs} (\texttt{/getPart}, \texttt{/createOrder}).
        \emph{Breaks:} uniform interface --- the URI vocabulary duplicates the
        method vocabulary, so nothing is safe, idempotent, or cacheable by
        inspection. \emph{Fix:} noun for the thing, method for the operation;
        reify true actions as resources (\texttt{/stock-adjustments}).
  \item \textbf{Tunneling everything through POST.} \emph{Breaks:} safety,
        idempotency, cacheability, visibility --- intermediaries must assume
        the worst about every request. \emph{Fix:} Level 2 method discipline;
        \texttt{GET} for reads so caches and conditional requests can work.
  \item \textbf{Returning 200 for errors} (error objects in a 200 body).
        \emph{Breaks:} self-descriptive messages --- monitors, retry loops, and
        caches all read status codes; a 200-wrapped failure may be
        \emph{cached}. \emph{Fix:} the Table~\ref{tab:ch3-status} lifecycle,
        with \texttt{application/problem+json} bodies for detail.
  \item \textbf{Leaking the database schema} (ORM objects serialized directly).
        \emph{Breaks:} independent evolvability --- schema migrations become
        breaking API changes. \emph{Fix:} an anti-corruption layer between
        storage and representation (Section~\ref{sec:ch3-representations}).
  \item \textbf{Inconsistent naming} (mixed singular/plural, camelCase beside
        kebab-case, \texttt{/order} beside \texttt{/warehouses}).
        \emph{Breaks:} the uniform interface's simplicity property --- every
        consumer now carries a lookup table of your exceptions. \emph{Fix:} one
        written grammar; lint it in CI (Listing~\ref{lst:ch3-urilint}).
  \item \textbf{Ignoring caching semantics} (no \texttt{Cache-Control}, no
        validators). \emph{Breaks:} cacheability --- intermediaries guess, and
        guessed caching is either absent (wasted scale) or wrong (stale data
        sold twice). \emph{Fix:} validators on every GET representation;
        explicit freshness policy, even if the policy is
        \texttt{no-cache} \cite{rfc9110}.
  \item \textbf{Deep nesting beyond two levels}
        (\texttt{/warehouses/WH-01/zones/Z-3/bins/B-17/parts/NX-1001}).
        \emph{Breaks:} evolvability --- the path ossifies one navigational
        route through a many-pathed graph; authorization and cache keys
        multiply. \emph{Fix:} top-level collections with query filters;
        containment paths only where lifetime truly depends on the parent.
  \item \textbf{Idempotency theater} (\texttt{PUT} that increments a counter).
        \emph{Breaks:} PUT's contract; safe retry becomes corruption.
        \emph{Fix:} if the operation is intrinsically non-idempotent
        (increment, append), use \texttt{POST} to a subordinate collection.
  \item \textbf{Versioned filenames-as-APIs} (\texttt{/api/v2/final/REAL/}).
        \emph{Breaks:} URI opacity and cool-URI discipline; version policy
        belongs in the contract, not in path folklore (Chapter 9).
  \item \textbf{Silent query-parameter divergence} (\texttt{?format=json}
        beside \texttt{Accept}). \emph{Breaks:} self-description and cache
        correctness (the URI key now varies by parameter a cache may ignore).
        \emph{Fix:} negotiate via headers with \texttt{Vary}, or make the
        variant a distinct URI deliberately --- never both.
\end{enumerate}

\begin{warningbox}
The most expensive anti-pattern is not on this list because it is meta:
\emph{labeling a Level 0/1 service ``REST''} in documentation. Consumers then
assume safety, idempotent retries, and cacheability that the service does not
provide, and build those assumptions into production traffic. Name your
architecture honestly; an ``HTTP RPC API'' is a respectable, supportable thing.
\end{warningbox}

%% ---------------------------------------------------------------------------
\section{Hands-On Production Lab \& Real-World Case Study}
\label{sec:ch3-lab}

\subsection{Lab Part A: Design the Warehouse Domain as REST Resources}

Work in the Northwind Robotics universe (all data synthetic). Target: 60--90
minutes, paper or a scratch OpenAPI file; no implementation required.

\begin{enumerate}
  \item \textbf{Inventory the domain.} Northwind's warehouse system tracks:
        warehouses; the parts catalog (Listing~\ref{lst:ch3-catalog}); stock
        levels (how many of a part sit in a warehouse); suppliers; purchase
        orders sent to suppliers; and inbound shipments against those orders.
  \item \textbf{Classify every concept} into the four archetypes
        (Section~\ref{sec:ch3-resources}): collections, documents, stores,
        controllers. Expect zero controllers; if you invented one, try to
        reify it.
  \item \textbf{Draw the resource graph.} Decide which relations are
        \emph{containment} (path hierarchy) versus \emph{association}
        (query filter or link). Compare with Figure~\ref{fig:ch3-resources}.
        Specifically: is \texttt{/warehouses/\{wid\}/stock/\{sku\}} justified?
        (Yes --- a stock level's identity is the pair (warehouse, part); it
        cannot exist without both.) Is \texttt{/suppliers/\{sid\}/purchase-orders/\{poid\}}
        justified? (No --- orders outlive any single navigational path; make
        \texttt{/purchase-orders?supplier=SID} instead.)
  \item \textbf{Write the URI table} for at least eight routes and lint it with
        \texttt{python uri\_lint.py yourtable.json} until clean.
  \item \textbf{Assign the uniform interface.} For each route, tabulate allowed
        methods, success codes, the three most likely failure codes, and the
        caching policy (validator, freshness) for each GET.
  \item \textbf{Design one hypermedia transition.} Purchase orders move
        \texttt{draft} $\rightarrow$ \texttt{submitted} $\rightarrow$
        \texttt{received} $\rightarrow$ \texttt{closed}. Express the legal
        transitions as state-dependent links on the order representation, as
        Listing~\ref{lst:ch3-hal} does for \texttt{restock}.
\end{enumerate}

\textbf{Exit criteria:} a lint-clean URI table, a completed method/status/cache
matrix, and a defensible two-sentence justification for every nested route.

\subsection{Lab Part B: Constraint Critique of a Deliberately Bad API}

The following route table ships from a fictional vendor, ``RoboDirect'' (any
resemblance to production APIs you have suffered is coincidental). Audit it
against every constraint and rule in this chapter before reading the key.

\begin{minted}{text}
POST   /api/getPartDetails            -> 200 always; errors in body
POST   /api/getPartDetails?id=NX-1001 -> same endpoint, query id
POST   /api/updateStockLevel          -> body: {partId, warehouseId, qty}
GET    /api/Parts/ListAll             -> 200, no ETag, no Cache-Control
POST   /api/warehouses/WH-01/zones/Z3/bins/B17/getContents
DELETE /api/purchaseOrder/delete/PO-2026-0417 -> 200 {"ok":true}
GET    /api/v2_final/parts/NX-1001.json
\end{minted}

\textbf{Answer key (abridged).} Every route tunnels through POST or embeds a
verb (R001), so nothing is safe, idempotent, or cacheable; errors hidden in 200
bodies defeat monitors and caches; \texttt{ListAll} returns an uncacheable,
unbounded collection with no validators; the five-segment warehouse path (R002)
ossifies one navigation path; \texttt{delete} in a DELETE path doubles the verb;
\texttt{v2\_final} and the \texttt{.json} suffix hard-code version and format
into the URI, defeating content negotiation and cool-URI discipline; and the
mixed \texttt{Parts}/\texttt{parts}/\texttt{purchaseOrder} naming fails R003
three ways. Estimated Richardson level: 0, trending 1. Rewrite left as
Self-Assessment Exercise 7.1.

\subsection{Case Study: The Redesign That Broke the Mobile Fleet}

\textbf{Context.} Northwind Robotics' field-technician app (a deployed fleet of
rugged Android devices scanning warehouse stock) consumed
\texttt{GET /api/inventory/\{sku\}}, which returned stock across all warehouses.
A platform team, having read a book much like this one, redesigned the API
``properly'': \texttt{/api/inventory/\{sku\}} became
\texttt{/api/v2/stock?part=\{sku\}}, launched in six weeks, and the old route
was deleted the same day. Within hours, the entire device fleet returned 404s;
technicians reverted to paper; two warehouses missed shipping cutoffs.

\textbf{Failure analysis.} The new URI design was, in isolation, better. The
failure was a discipline failure, not a design failure. (1) \emph{URIs are
contracts}: clients had hard-coded paths (violating opacity, as clients
inevitably do), so renaming was a breaking change regardless of how ``clean''
the new grammar was. (2) \emph{No coexistence window}: mobile fleets upgrade on
app-store timescales of months; the old route must live until measured traffic
reaches zero, returning deprecation signals (\texttt{Deprecation} and
\texttt{Sunset} headers per RFC 8594 and RFC 9745 practice). (3) \emph{No
consumer observability}: nobody could answer ``who still calls the old route?''
because access logs were never keyed by client version. (4) \emph{Versioning as
renaming}: \texttt{v2} was minted for a change that a new \emph{resource}
(\texttt{/stock-levels}) alongside the old one would have made non-breaking.

\textbf{Recovery and doctrine.} The old route was restored behind a 308 redirect
to a compatibility shim; a deprecation schedule with dated sunset headers was
published; per-client-version request logging was added at the gateway; and the
governance rule adopted was: \emph{additive changes ship immediately; breaking
changes ship only behind measured coexistence}. Chapters 9 and 20 turn this war
story into process.

%% ---------------------------------------------------------------------------
\section{Self-Assessment Exercises \& Deep-Dive Questions}
\label{sec:ch3-assessment}

\begin{enumerate}
  \item \textbf{(7.1) Rewrite RoboDirect.} Produce a lint-clean route table for
        the Lab Part B API, with a method/status/cache matrix for every route.
  \item \textbf{(7.2)} State statelessness precisely enough to decide: is a
        request carrying a session cookie stateless? A bearer token? An OAuth
        access token validated against a token-introspection endpoint? Defend
        each answer from the definition.
  \item \textbf{(7.3)} Northwind sells ``today's promotions'' --- a resource
        whose value changes daily. Design its URI, freshness headers, and
        validators. What breaks if a shared cache serves it eight hours stale?
  \item \textbf{(7.4)} Extend Listing~\ref{lst:ch3-catalog} with
        \texttt{HEAD} correctness verified by the selftest, and with a
        \texttt{POST /parts}-side \texttt{Idempotency-Key} header. Which
        constraint does each addition strengthen?
  \item \textbf{(7.5)} Implement an RFC 6902 JSON Patch evaluator (operations
        \texttt{add}, \texttt{remove}, \texttt{replace}, \texttt{test}) and wire
        it into the catalog's PATCH alongside Merge Patch. Which media type
        selects which evaluator?
  \item \textbf{(7.6) Deep dive.} Fielding calls the uniform interface
        trade-off ``efficiency versus evolvability.'' Construct one concrete
        Northwind scenario where the uniform interface is measurably less
        efficient than a bespoke RPC, and state what is purchased with that
        inefficiency.
  \item \textbf{(7.7) Deep dive.} Argue both sides: should
        \texttt{/warehouses/\{wid\}/stock/\{sku\}} expose \texttt{DELETE}, and
        if so, what does it mean --- zero stock, or no record? What do 404 and
        410 each communicate here?
  \item \textbf{(7.8) Deep dive.} Your organization mandates JSON:API. Which
        decisions in this chapter does it make for you, which does it forbid,
        and where does its opinion cost you flexibility the catalog used?
\end{enumerate}

%% ---------------------------------------------------------------------------
\section{Interview Q\&A Vault}
\label{sec:ch3-qa}

Thirty-six questions, ordered junior $\rightarrow$ staff. Answers are calibrated
to what a strong candidate says in two to six sentences.

\subsection{Junior (Q1--Q10)}

\begin{enumerate}[label=\textbf{Q\arabic*.}, leftmargin=3.2em, itemsep=0.6em]
  \item \textbf{Is REST a protocol or a standard?}\\
        \emph{Answer:} Neither. REST is an architectural \emph{style} --- a set
        of constraints defined in Fielding's dissertation
        \cite{fielding2000rest} --- that induces properties like scalability
        and evolvability. HTTP is a protocol; OpenAPI is a description
        standard; REST constrains how you use them.
  \item \textbf{State Fielding's six constraints.}\\
        \emph{Answer:} Client--server, statelessness, cacheability, uniform
        interface, layered system, and (optional) code-on-demand. Bonus depth:
        the uniform interface decomposes into resource identification,
        manipulation through representations, self-descriptive messages, and
        HATEOAS.
  \item \textbf{What properties does statelessness induce?}\\
        \emph{Answer:} Visibility (any request is understandable in isolation),
        reliability (requests can simply be retried after failure), and
        scalability (no session affinity; any replica serves any request).
        Cost: session data travels on every request.
  \item \textbf{What is a resource?}\\
        \emph{Answer:} Any named conceptual target of a hypertext reference ---
        a document, a collection, a service, even ``today's price of NX-1001.''
        Clients never touch the resource, only representations of its state.
  \item \textbf{Name the four resource archetypes.}\\
        \emph{Answer:} Document (singular object), collection (server-managed
        directory), store (client-managed entry repository), and controller
        (verb-like escape hatch, to be minimized by reifying actions as
        resources).
  \item \textbf{Why ``nouns, not verbs'' in URIs?}\\
        \emph{Answer:} The HTTP method already supplies the verb. Verbs in URIs
        duplicate that vocabulary with unbounded private words, destroying the
        safety, idempotency, and cacheability signals that generic
        infrastructure depends on.
  \item \textbf{What do safe and idempotent mean? Give one example each.}\\
        \emph{Answer:} Safe: read-only, no accountability for state change
        (GET). Idempotent: repeating the request has the same effect as one
        (PUT, DELETE). Note idempotency is about server \emph{state}, not the
        response --- a repeated DELETE returns 404 but is still idempotent.
  \item \textbf{What status code should a successful POST that creates a
        resource return, and with which header?}\\
        \emph{Answer:} 201 Created with a \texttt{Location} header pointing at
        the new resource's URI; ideally also its representation and ETag.
  \item \textbf{What is the difference between 400 and 422?}\\
        \emph{Answer:} 400: the request is malformed --- unparseable JSON,
        missing required fields, wrong types. 422: syntactically fine but
        semantically unprocessable --- a well-formed request whose values
        violate business rules (e.g., negative stock).
  \item \textbf{What does a 405 response require?}\\
        \emph{Answer:} An \texttt{Allow} header listing the methods the
        resource does support, so clients can self-correct without
        documentation archaeology.
\end{enumerate}

\subsection{Mid-Level (Q11--Q20)}

\begin{enumerate}[label=\textbf{Q\arabic*.}, leftmargin=3.2em, itemsep=0.6em, resume]
  \item \textbf{Explain the Richardson Maturity Model.}\\
        \emph{Answer:} Four levels diagnosing which web technologies an HTTP API
        exploits \cite{richardson2013restful}: L0 tunnels RPC through one URI
        and POST; L1 introduces resource URIs; L2 uses methods and status codes
        with correct semantics; L3 adds hypermedia controls (HATEOAS). It is
        diagnostic, not a moral ladder.
  \item \textbf{PUT versus PATCH?}\\
        \emph{Answer:} PUT is full replacement --- the body becomes the
        resource's new state; omitted fields are reset; idempotent by
        construction. PATCH is partial modification whose semantics come from
        the patch media type --- JSON Merge Patch (RFC 7386: merge, null
        deletes) or JSON Patch (RFC 6902: ordered atomic operations with a
        \texttt{test} precondition).
  \item \textbf{When is a repeated DELETE not an error?}\\
        \emph{Answer:} Always, by design. The first DELETE returns 204; the
        second returns 404. Both leave the server in the same state, which is
        exactly why DELETE is idempotent and safe for retry loops --- the
        client should treat 404-after-DELETE as confirmation, not failure.
  \item \textbf{How do ETags and conditional requests work?}\\
        \emph{Answer:} The server attaches a validator (\texttt{ETag}) to a
        representation. A client revalidating sends \texttt{If-None-Match}; a
        match yields 304 with no body (cache efficiency). A client writing
        sends \texttt{If-Match}; a mismatch yields 412 (optimistic concurrency
        --- no lost updates) \cite{rfc9110}.
  \item \textbf{Design a URI scheme for: parts, stock per warehouse, purchase
        orders to suppliers.}\\
        \emph{Answer:} \texttt{/parts}, \texttt{/parts/\{sku\}};
        \texttt{/warehouses/\{wid\}/stock/\{sku\}} (containment: stock's
        identity is the pair); \texttt{/purchase-orders} top-level with
        \texttt{?supplier=} and \texttt{?warehouse=} filters because orders
        outlive any single navigation path. Depth $\le$ 2, kebab-case, plural
        collections.
  \item \textbf{Query string or path segment --- how do you decide?}\\
        \emph{Answer:} Path identifies \emph{which} resource; query modulates
        \emph{how} it is returned: filtering, sorting, pagination, projection.
        If deleting the parameter addresses a different thing, it belongs in
        the path.
  \item \textbf{Singular or plural collection names?}\\
        \emph{Answer:} The debate is noise; consistency is the law. Plural wins
        on three arguments: collection and member URIs share a stem
        (\texttt{/parts}, \texttt{/parts/NX-1001}); it reads correctly for both
        many and one-of-many; and singletons stay expressible as documents.
        Mixing is the only indefensible answer.
  \item \textbf{Are \texttt{/parts} and \texttt{/parts/} the same resource?}\\
        \emph{Answer:} No --- URIs are compared character-wise, so they are
        distinct identifiers with distinct cache keys. Publish the bare form
        and answer the slashed form with a cacheable 301 redirect; never alias
        silently.
  \item \textbf{What is URI opacity and what does it buy?}\\
        \emph{Answer:} Clients treat URIs as opaque: store, compare,
        dereference --- never parse or construct. It buys the server freedom to
        restructure (gateways, sharding, renaming) without breaking consumers,
        and it is the precondition for hypermedia-driven evolution.
  \item \textbf{When would you return 410 instead of 404?}\\
        \emph{Answer:} 404 says ``not here, maybe never was, maybe
        elsewhere''; 410 says ``was here, deliberately retired, do not ask
        again.'' Use 410 for sunset resources so crawlers, caches, and clients
        can permanently drop references.
\end{enumerate}

\subsection{Senior (Q21--Q28)}

\begin{enumerate}[label=\textbf{Q\arabic*.}, leftmargin=3.2em, itemsep=0.6em, resume]
  \item \textbf{What does the layered-system constraint permit, and what does
        it cost?}\\
        \emph{Answer:} It permits intermediaries --- gateways, shared caches,
        CDNs --- that components cannot see past, enabling load balancing,
        edge security, and legacy encapsulation. It costs per-hop latency,
        which the cache constraint is expected to amortize
        \cite{fielding2000rest}.
  \item \textbf{Why is code-on-demand optional, and where do you actually see
        it?}\\
        \emph{Answer:} It reduces visibility --- intermediaries cannot reason
        about downloaded behavior --- so Fielding brackets it as the one
        optional constraint. Its canonical realization is JavaScript to
        browsers; machine-to-machine APIs almost always omit it.
  \item \textbf{When does hypermedia pay off, and when is it over-engineering?}\\
        \emph{Answer:} It pays when consumers outnumber your coordination
        bandwidth: long-lived public APIs, server-driven workflows whose legal
        transitions vary per state, generic clients. It is over-engineering for
        first-party clients shipping in lockstep, where static contracts
        (OpenAPI) plus selective links (pagination, state transitions) capture
        most of the value \cite{richardson2013restful}.
  \item \textbf{Compare HAL, JSON:API, and Collection+JSON.}\\
        \emph{Answer:} HAL is minimal --- \texttt{\_links},
        \texttt{\_embedded}, curies; reads only. JSON:API is a full protocol:
        resource objects, relationships, compound documents, pagination and
        error conventions, write semantics. Collection+JSON is the purest
        hypermedia: the representation carries queries and a write template,
        enabling fully generic clients.
  \item \textbf{How do you keep an API from becoming a serialized database?}\\
        \emph{Answer:} An anti-corruption layer: representations designed for
        consumers (structured money, curated fields), storage designed for
        integrity, explicit translation between. Signals you have failed:
        column names in JSON, ORM objects passed to the encoder, schema
        migrations breaking clients \cite{kleppmann2017ddia}.
  \item \textbf{Design PATCH for a resource edited concurrently by many
        clients.}\\
        \emph{Answer:} JSON Patch (RFC 6902) with \texttt{test} operations as
        preconditions and \texttt{If-Match} on the ETag; apply atomically or
        not at all; return 409/412 with a problem detail so the client re-reads
        and re-derives its patch. Merge Patch is the wrong tool --- no
        atomicity, no preconditions.
  \item \textbf{A stakeholder demands \texttt{?format=csv} instead of content
        negotiation. Trade-offs?}\\
        \emph{Answer:} It works and is browser-testable, but it forks the
        identifier space (two URIs, one resource), complicates cache keys
        unless \texttt{Vary}-equivalent discipline is applied to the query, and
        bypasses q-values. Acceptable as a documented alias; the
        \texttt{Accept} header with \texttt{Vary: Accept} remains the canonical
        mechanism \cite{rfc9110}.
  \item \textbf{Your API must expose a batch ``revalue inventory'' operation.
        RESTful design?}\\
        \emph{Answer:} Reify the operation: \texttt{POST /inventory-revaluations}
        creates a revaluation resource with its own URI, status, and audit
        trail; \texttt{GET} on it reports progress. You get idempotent
        submission (via keys), observable lifecycle, and cancelability ---
        none of which a \texttt{POST /revalueInventory} verb-endpoint can
        offer.
\end{enumerate}

\subsection{Staff (Q29--Q36)}

\begin{enumerate}[label=\textbf{Q\arabic*.}, leftmargin=3.2em, itemsep=0.6em, resume]
  \item \textbf{``REST doesn't scale to our read patterns --- clients need
        aggregates from five resources.'' Respond.}\\
        \emph{Answer:} REST never forbids server-composed resources: a
        \texttt{/dashboards/warehouse-summary} resource is a first-class,
        cacheable document. What REST forbids is letting clients dictate
        per-request join graphs ad hoc. Compose on the server behind a stable
        resource, or admit you need a query paradigm (GraphQL, Chapter 12) ---
        the honest trade, not a tunnelled query POST pretending to be REST.
  \item \textbf{How do you evolve a Level 2 API for five years without
        \texttt{/v2}?}\\
        \emph{Answer:} Additive-only changes, tolerant readers, never reuse
        retired field names, closed-by-default schemas in review
        (\texttt{additionalProperties: false} as a forcing function),
        deprecation headers with measured coexistence, and consumer-keyed
        access logging so every removal decision is data-driven. Version
        banners become a rare, deliberate event (Chapter 9)
        \cite{gehl2021apidesign}.
  \item \textbf{Define a governance model for URI and contract consistency
        across forty teams.}\\
        \emph{Answer:} A short written style guide (archetypes, grammar, status
        matrix), executable linting in CI with stable rule IDs (the
        Listing~\ref{lst:ch3-urilint} pattern), a review gate only for
        cross-cutting changes, and published exception/waiver mechanics.
        Consistency enforced by tooling scales; consistency enforced by
        meetings does not \cite{gehl2021apidesign}.
  \item \textbf{Statelessness versus connection-oriented protocols: when do you
        abandon the constraint?}\\
        \emph{Answer:} When the interaction \emph{is} a session: streams,
        collaborative editing, subscriptions. The disciplined move is to
        declare the boundary --- REST for command/control resources, WebSocket
        or gRPC streaming (Chapters 11, 13) for the conversational edge ---
        rather than smuggling session state into ``REST'' endpoints.
  \item \textbf{A cache served stale stock levels and inventory was oversold.
        Root-cause through the constraints.}\\
        \emph{Answer:} The cacheability constraint was violated by omission:
        responses carried no freshness metadata, so an intermediary guessed.
        Fix at the origin: validators plus explicit \texttt{Cache-Control}
        (\texttt{no-cache} revalidation for stock), and treat any response
        without a caching policy as a defect in review \cite{rfc9110}.
  \item \textbf{When is RPC legitimately preferable to REST? Say so without
        hedging.}\\
        \emph{Answer:} Low-latency internal service-to-service calls with
        schema-first contracts and no intermediaries (gRPC's home); operations
        that are intrinsically computations over transient inputs
        (\texttt{calculateRoute}); and teams whose consumers all ship in
        lockstep. REST buys evolvability and interoperability with the web's
        infrastructure; pay for it only when you need it.
  \item \textbf{Design the migration from a Level 0 SOAP estate to resource
        orientation without a big bang.}\\
        \emph{Answer:} Strangle at the edge: stand up a gateway translating the
        ten highest-traffic RPC operations into resources; publish OpenAPI for
        the new surface; dual-run with reconciliation dashboards; migrate
        consumers by cohort with sunset headers on the legacy endpoint; never
        rename and restructure in the same release (the Lab case study).
  \item \textbf{What would you instrument to prove REST's claimed properties
        are actually materializing?}\\
        \emph{Answer:} Cache hit ratio and revalidation rate (cacheability);
        cross-replica request distribution without affinity (statelessness);
        share of 4xx/5xx correctly classified by status code alone
        (visibility/self-description); consumer breakage rate per contract
        change (evolvability); and p99 added latency per intermediary layer
        (the layered system's bill coming due).
\end{enumerate}

%% ---------------------------------------------------------------------------
\section{External Bibliography \& Further Reading}
\label{sec:ch3-bibliography}

\begin{itemize}
  \item \textbf{Roy T. Fielding,} \emph{Architectural Styles and the Design of
        Network-based Software Architectures,} Ph.D. dissertation, UC Irvine,
        2000 \cite{fielding2000rest}. Chapter 5 is the primary source for
        everything in Section~\ref{sec:ch3-constraints}; read it before any
        secondary source, including this book.
  \item \textbf{Leonard Richardson, Mike Amundsen, Sam Ruby,}
        \emph{RESTful Web APIs,} O'Reilly, 2013 \cite{richardson2013restful}.
        The maturity model, Collection+JSON, and the most honest treatment of
        hypermedia economics in print.
  \item \textbf{Mark Mass\'e,} \emph{REST API Design Rulebook,} O'Reilly, 2011.
        The WRML-influenced four-archetype taxonomy (document, collection,
        store, controller) and URI grammar conventions used in
        Section~\ref{sec:ch3-resources}.
  \item \textbf{JJ Geewax,} \emph{API Design Patterns,} Manning, 2021
        \cite{gehl2021apidesign}. A pattern catalogue for resource layout,
        pagination, long-running operations, and versioning; the vocabulary of
        Section~\ref{sec:ch3-product}'s governance discussion.
  \item \textbf{RFC 9110,} \emph{HTTP Semantics,} IETF, 2022 \cite{rfc9110}.
        The normative source for method properties, status codes, content
        negotiation, and validators; cite it in design reviews, not folklore.
  \item \textbf{RFC 8288,} \emph{Web Linking,} IETF, 2017 \cite{rfc8288}. Typed
        links and relation types --- the mechanical foundation of HATEOAS.
  \item \textbf{RFC 7386} (\emph{JSON Merge Patch}) and \textbf{RFC 6902}
        (\emph{JSON Patch}), IETF. The two patch grammars compared in
        Section~\ref{sec:ch3-uniform}; both short enough to read in one
        sitting.
  \item \textbf{OpenAPI Specification 3.1,} OpenAPI Initiative
        \cite{openapi31}. The contract format for the design-first workflow;
        note that 3.1 aligns fully with JSON Schema 2020-12.
  \item \textbf{JSON:API specification} (jsonapi.org). Even if you never adopt
        it, its pagination, error-object, and compound-document conventions are
        a free design review of your own conventions.
  \item \textbf{Zalando RESTful API Guidelines} and \textbf{Microsoft REST API
        Guidelines.} Two battle-tested corporate style guides; read them to see
        how the debates of Section~\ref{sec:ch3-uri-grammar} get settled in
        practice (kebab-case, plural collections, explicit null policy).
  \item \textbf{Martin Kleppmann,} \emph{Designing Data-Intensive Applications,}
        O'Reilly, 2017 \cite{kleppmann2017ddia}. Chapter on encoding and
        evolution: the theoretical basis for anti-corruption representations
        and tolerant readers.
  \item \textbf{The Twelve-Factor App} \cite{twelvefactor}. Process
        disposability and backing-services discipline: operational cousins of
        statelessness.
\end{itemize}

%% ---------------------------------------------------------------------------
\section{Capstone Projects}
\label{sec:ch3-capstones}

Five progressive projects. Each produces artifacts suitable for a design-review
portfolio. All are offline and use synthetic Northwind Robotics data only.

\subsection{Capstone 3.1 [junior] --- Model the Library Domain}

\textbf{Objective.} Practice pure resource modeling on an unfamiliar domain
before touching code.

\textbf{Inputs/Datasets.} A fictional circulating library: books (with editions
and copies), members, loans, holds, fines. Invent all data; at least ten books,
four members, three active loans.

\textbf{Steps.}
\begin{enumerate}
  \item Classify every concept into the four archetypes; reify at least two
        ``actions'' (returning a book, placing a hold) as resources.
  \item Draw the resource graph; justify every containment edge in one
        sentence.
  \item Write the full route table and lint it with
        Listing~\ref{lst:ch3-urilint} until clean.
  \item Produce the method/status/cache matrix for every route.
\end{enumerate}

\textbf{Acceptance Criteria.}
\begin{itemize}
  \item[$\square$] Zero lint findings (rules R001--R005).
  \item[$\square$] No controller archetype without a written reification
        failure analysis.
  \item[$\square$] Nesting depth never exceeds two.
  \item[$\square$] Every GET route has a stated validator and freshness policy.
\end{itemize}

\textbf{Stretch Goals.} Add sparse fieldsets and an \texttt{?embed=} policy;
write the JSON Merge Patch document that renews a loan.

\subsection{Capstone 3.2 [junior] --- Run and Extend the Stdlib Catalog}

\textbf{Objective.} Operate Listing~\ref{lst:ch3-catalog} and prove its
constraint compliance empirically.

\textbf{Inputs/Datasets.} \texttt{catalog\_api.py}; \texttt{curl.exe} (ships
with Windows 10+); PowerShell.

\textbf{Steps.}
\begin{enumerate}
  \item Run \texttt{python catalog\_api.py -{}-selftest} and map each assertion
        to the constraint it demonstrates.
  \item Serve on port 8080; with \texttt{curl -i}, demonstrate: 200 with ETag;
        304 via \texttt{If-None-Match}; 201 with \texttt{Location}; 409 on
        duplicate SKU; 412 via a stale \texttt{If-Match} on PUT; 406 via
        \texttt{Accept: application/xml}; CSV via \texttt{Accept: text/csv}.
  \item Add one category, \texttt{batteries}, end-to-end (validation set, seed
        data, one new selftest assertion).
\end{enumerate}

\textbf{Acceptance Criteria.}
\begin{itemize}
  \item[$\square$] All seven demonstrations captured as transcripts in a
        \texttt{lab-notes.md}.
  \item[$\square$] Selftest passes after your extension.
  \item[$\square$] A one-paragraph memo: which constraint would break first if
        you added server-side login sessions, and why.
\end{itemize}

\textbf{Stretch Goals.} Implement \texttt{HEAD}-only verification; add
\texttt{Sunset}/\texttt{Deprecation} headers to a deliberately retired part.

\subsection{Capstone 3.3 [mid] --- Productionize the URI Linter}

\textbf{Objective.} Turn Listing~\ref{lst:ch3-urilint} into a governance tool a
real platform team could run in CI.

\textbf{Inputs/Datasets.} \texttt{uri\_lint.py}; three route tables you author:
one clean, one failing each rule, one extracted by hand from a public API's
documentation (retyped, synthetic base URLs).

\textbf{Steps.}
\begin{enumerate}
  \item Add a waiver mechanism (\texttt{\# lint-waive: R002 reason} in the
        table's metadata) with expiry dates.
  \item Add rules: R006 (identifier segments must pair with a collection),
        R007 (query parameters documented in the table must be kebab-case).
  \item Emit machine-readable output (\texttt{-{}-format json}) alongside the
        text report.
  \item Write \texttt{pytest}-style unit tests (stdlib \texttt{unittest} is
        acceptable) covering every rule, including false-positive cases.
\end{enumerate}

\textbf{Acceptance Criteria.}
\begin{itemize}
  \item[$\square$] Exit codes 0/1/2 behave as documented.
  \item[$\square$] Waivers suppress only their cited rule and path, and expire.
  \item[$\square$] Full rule test coverage, all green and deterministic.
  \item[$\square$] JSON report validates against a schema you write.
\end{itemize}

\textbf{Stretch Goals.} Accept a real OpenAPI 3.1 document as input; diff two
route tables and classify each delta as breaking or additive.

\subsection{Capstone 3.4 [senior] --- Port an RPC API to Resource Orientation}

\textbf{Objective.} Execute the Section~\ref{sec:ch3-lab} migration doctrine on
a realistic legacy surface.

\textbf{Inputs/Datasets.} The RoboDirect API from Lab Part B, extended by you to
twelve endpoints (add shipping, invoicing, and technician scheduling RPCs); a
synthetic two-week access log you generate showing per-endpoint traffic.

\textbf{Steps.}
\begin{enumerate}
  \item Model the full RoboDirect domain as resources; produce the target
        route table and lint it.
  \item Write the OpenAPI 3.1 contract for the new surface, design-first.
  \item Implement compatibility shims (301/308 redirects or proxy routes)
        mapping every legacy call to a new resource, with
        \texttt{Deprecation} and \texttt{Sunset} headers.
  \item Write a migration plan keyed by the access log: consumer cohorts,
        coexistence windows, measurable exit criteria per cohort.
\end{enumerate}

\textbf{Acceptance Criteria.}
\begin{itemize}
  \item[$\square$] Every legacy endpoint has a mapped resource or a written
        reification failure analysis.
  \item[$\square$] Contract passes Listing~\ref{lst:ch3-urilint} with zero
        findings.
  \item[$\square$] Shims demonstrably preserve behavior for the twelve legacy
        calls (before/after transcripts).
  \item[$\square$] Migration plan names dates, owners, and kill criteria.
\end{itemize}

\textbf{Stretch Goals.} Add hypermedia transition links to the two
workflow-heavy resources (shipments, invoices); quantify payload and
round-trip deltas versus the RPC originals.

\subsection{Capstone 3.5 [staff] --- Design-Review Report on a Public API}

\textbf{Objective.} Produce a staff-grade architectural critique of a real,
publicly documented API (documentation review only --- no production traffic;
a captured, synthetic replay of documented examples is the dataset).

\textbf{Inputs/Datasets.} The public documentation and OpenAPI document of one
large API of your choice; this chapter's constraint tables; the
Listing~\ref{lst:ch3-urilint} linter run against its published paths.

\textbf{Steps.}
\begin{enumerate}
  \item Score the API against all six constraints with evidence (captured
        responses, documented behaviors).
  \item Place it on the Richardson scale per subsystem --- large APIs are
        rarely uniform.
  \item Quantify: lint findings by rule; status-code usage distribution;
        cacheability of GET responses; hypermedia coverage.
  \item Write a prioritized remediation backlog with effort/impact estimates
        and a coexistence-safe migration approach for the top three items.
\end{enumerate}

\textbf{Acceptance Criteria.}
\begin{itemize}
  \item[$\square$] Every claim cites observable evidence (header capture,
        documented example, lint output).
  \item[$\square$] At least one finding credits the API for a sophisticated
        choice --- pure hatchet jobs fail review.
  \item[$\square$] Remediation backlog survives the case-study doctrine: no
        rename-and-restructure in a single step.
  \item[$\square$] Report is eight pages or fewer and readable by a
        non-specialist executive.
\end{itemize}

\textbf{Stretch Goals.} Repeat the review five years apart using archived
documentation; measure whether the vendor's evolvability claims materialized.

%% ---------------------------------------------------------------------------
\section{Python Dependencies Appendix}
\label{sec:ch3-dependencies}

The chapter is stdlib-first: Listings~\ref{lst:ch3-catalog} and
\ref{lst:ch3-hal} run on a bare Python 3.11 installation. Only the route-table
validator has an optional third-party need.

\begin{minted}{text}
# requirements.txt -- Chapter 3 (Python 3.11+)
# Core listings (catalog API, HAL serializer): standard library only.
pyyaml>=6.0        # uri_lint.py: YAML route tables (JSON input works without it)
jsonschema>=4.21   # optional: structurally validate OpenAPI 3.1 documents
\end{minted}

\subsection{pyyaml ($\geq$ 6.0)}

\textbf{Description.} The de-facto YAML 1.1/1.2 parser for Python; used by
\texttt{uri\_lint.py} only when the route table is authored in YAML (the native
format of most OpenAPI documents).

\textbf{Why included.} YAML is the authoring format of the API-design ecosystem;
rejecting it would make the linter useless against real contracts. It is a
mature, widely pinned dependency with no runtime services.

\textbf{Alternatives.} \texttt{ruamel.yaml} (round-trip preservation; heavier);
the standard library alone (accept JSON only --- the linter already does).

\textbf{Installation (PowerShell).}
\begin{minted}{bash}
python -m pip install "pyyaml>=6.0"
\end{minted}

\subsection{jsonschema ($\geq$ 4.21)}

\textbf{Description.} A JSON Schema validator implementing the 2020-12 draft
that OpenAPI 3.1 adopts \cite{openapi31}.

\textbf{Why included.} Optional belt-and-braces: validate that an OpenAPI
document's schemas are well-formed and that example payloads conform, before
the contract reaches consumers. Not required by any listing.

\textbf{Alternatives.} \texttt{openapi-spec-validator} (validates the OpenAPI
document itself, built on jsonschema); \texttt{check-jsonschema} (CLI wrapper).

\textbf{Installation (PowerShell).}
\begin{minted}{bash}
python -m pip install "jsonschema>=4.21"
\end{minted}

%% ---------------------------------------------------------------------------
\section{Chapter Summary \& Key Takeaways}

\begin{itemize}
  \item \textbf{REST is a style, not a protocol.} Six constraints ---
        client--server, statelessness, cacheability, uniform interface, layered
        system, optional code-on-demand --- are traded for scalability,
        visibility, reliability, and evolvability \cite{fielding2000rest}.
  \item \textbf{Diagnose before you judge.} The Richardson Maturity Model
        (POX $\rightarrow$ resources $\rightarrow$ verbs $\rightarrow$
        hypermedia) tells you which constraints an API actually honors
        \cite{richardson2013restful}. ``JSON over HTTP'' tells you nothing.
  \item \textbf{Model things, not operations.} Nouns in URIs, the four
        archetypes, reified actions, containment-only nesting (depth $\le 2$),
        kebab-case, plural collections, bare trailing-slash policy, path for
        identity and query for modulation.
  \item \textbf{Let the uniform interface work.} Safety and idempotency are
        contracts your retries and caches depend on; status codes are the
        server's half of the conversation; ETags plus conditional requests give
        caching and optimistic concurrency for free \cite{rfc9110}.
  \item \textbf{PATCH means choosing a patch grammar.} Merge Patch for
        forgiving updates, JSON Patch for atomic, preconditioned, concurrent
        editing.
  \item \textbf{Hypermedia is an investment decision.} Links and affordances
        \cite{rfc8288} buy restructuring freedom; adopt selectively unless your
        consumers genuinely outnumber your coordination bandwidth.
  \item \textbf{Representations are products, not dumps.} Anti-corruption
        layering, additive evolution, tolerant readers, and never reusing field
        names keep APIs evolving for years \cite{kleppmann2017ddia,gehl2021apidesign}.
  \item \textbf{Govern with code.} A linted route table and a design-first
        OpenAPI contract \cite{openapi31} beat a wiki of good intentions.
\end{itemize}

%% ---------------------------------------------------------------------------
\section{Quick-Reference Card}

\begin{table}[h]
  \centering
  \small
  \begin{tabularx}{\textwidth}{l X}
    \toprule
    \textbf{Constraint} & \textbf{One-line test} \\
    \midrule
    Client--server & Can the UI and the datastore each change without the other? \\
    Statelessness & Is every request self-contained --- no server session context? \\
    Cacheability & Does every response say whether and how it may be cached? \\
    Uniform interface & Resources identified; manipulated via representations; messages self-describing; hypermedia drives state? \\
    Layered system & Could an intermediary be inserted without semantic change? \\
    Code-on-demand (opt.) & If scripts are served, is the visibility loss accepted in writing? \\
    \bottomrule
  \end{tabularx}
  \caption{Six-constraint compliance checklist.}
\end{table}

\begin{table}[h]
  \centering
  \small
  \begin{tabular}{l c c l}
    \toprule
    \textbf{Method} & \textbf{Safe} & \textbf{Idempotent} & \textbf{Success} \\
    \midrule
    GET/HEAD & yes & yes & 200 / 304 \\
    POST     & no  & no  & 201 (+ Location) or 200 \\
    PUT      & no  & yes & 200 / 204 / 201 (412 on stale If-Match) \\
    PATCH    & no  & (design for it) & 200 / 204 \\
    DELETE   & no  & yes & 204 (404 on repeat is fine) \\
    OPTIONS  & yes & yes & 204 + Allow \\
    \bottomrule
  \end{tabular}
  \caption{Method semantics at a glance \cite{rfc9110}.}
\end{table}

\begin{table}[h]
  \centering
  \small
  \begin{tabularx}{\textwidth}{l X X}
    \toprule
    \textbf{RMM level} & \textbf{Signature} & \textbf{What you gain} \\
    \midrule
    0 --- POX & one URI, POST, 200-always & nothing; HTTP as tunnel \\
    1 --- Resources & many noun URIs, one verb & addressability \\
    2 --- Verbs & methods + status codes + validators & caching, safe retry, visibility \\
    3 --- Hypermedia & representations carry links/forms & restructuring freedom, runtime discovery \\
    \bottomrule
  \end{tabularx}
  \caption{Richardson Maturity Model diagnostic \cite{richardson2013restful}.}
\end{table}

\begin{minted}{text}
URI grammar: lowercase kebab-case . plural collections . depth <= 2 .
  no verbs . no trailing slash (301 the slashed form) .
  path = identity, query = modulation . URIs are opaque to clients
\end{minted}

\section{Standardized Evidence Addendum}
\subsection{Benchmark Snapshot Template}
\begin{table}[h]
  \centering
  \small
  \begin{tabularx}{\textwidth}{l c c c c}
    \toprule
    \textbf{Config} & \textbf{p50 latency} & \textbf{p99 latency} & \textbf{Error rate} & \textbf{Throughput} \\
    \midrule
    Baseline & -- & -- & -- & 1.00x \\
    Candidate A & -- & -- & -- & -- \\
    Candidate B & -- & -- & -- & -- \\
    \bottomrule
  \end{tabularx}
  \caption{Minimum benchmark table to keep per chapter.}
\end{table}

\subsection{Request Trace Template}
\begin{itemize}
  \item \textbf{Request:} one representative API call for this chapter's topic.
  \item \textbf{Before:} behavior (headers, payload, latency, failure mode) with the naive approach.
  \item \textbf{After:} behavior with the technique in this chapter.
  \item \textbf{Result:} what changed in correctness, resilience, latency, and operability.
\end{itemize}

\subsection{Cost and Latency Budget Snapshot}
\begin{table}[h]
  \centering
  \small
  \begin{tabularx}{\textwidth}{l c c}
    \toprule
    \textbf{Stage} & \textbf{Latency budget} & \textbf{Cost driver} \\
    \midrule
    TLS / connection setup & -- & handshake round trips, keep-alive hit rate \\
    Request validation & -- & schema complexity, CPU per request \\
    Business logic and I/O & -- & downstream fan-out, query cost \\
    Serialization and egress & -- & payload size, compression ratio \\
    \bottomrule
  \end{tabularx}
  \caption{Operational budget template for this chapter's technique.}
\end{table}

\section{Configuration Preset Template}
\begin{minted}{text}
# api_policy.yaml
policy_version: "v1.0.0"
component: "<chapter_topic>"
timeout_ms: 0
retry_budget: 0
fallback_strategy: "fail_fast"
rollback_trigger: "error_budget_burn_gt_threshold"
\end{minted}

\section{CI Quality Gate Specification}
\begin{itemize}
  \item contract tests green against the pinned OpenAPI/AsyncAPI/Protobuf spec,
  \item no breaking schema changes without a version bump,
  \item error responses conform to RFC 9457 problem-details shape,
  \item benchmark deltas within approved regression bounds (latency and throughput),
  \item policy version and rollback plan present in release metadata.
\end{itemize}

\section{Incident Runbook Template}
\begin{enumerate}
  \item Detect regression from SLO burn-rate alerts or consumer reports.
  \item Scope impacted endpoints, versions, and consumer segments.
  \item Contain impact (rollback, feature flag, rate-limit tightening, or traffic shift).
  \item Diagnose with traces, structured logs, and contract diffs.
  \item Recover with validated fix and verified rollout procedure.
  \item Prevent recurrence via new regression tests and CI gates.
\end{enumerate}

\section{Cross-Chapter Decision Card}
\begin{table}[h]
  \centering
  \small
  \begin{tabularx}{\textwidth}{l X X}
    \toprule
    \textbf{Dimension} & \textbf{Choose this approach when} & \textbf{Prefer an alternative when} \\
    \midrule
    Use case & -- & -- \\
    Latency budget & -- & -- \\
    Operational cost & -- & -- \\
    Team maturity & -- & -- \\
    \bottomrule
  \end{tabularx}
  \caption{Decision card to complete per chapter.}
\end{table}

\section{ROI Model and Build-vs-Buy}
\begin{itemize}
  \item \textbf{Build when:} the capability is differentiating, compliance forbids vendors, or scale makes vendor pricing irrational.
  \item \textbf{Buy when:} the capability is commodity (auth, gateways, observability), time-to-market dominates, or staffing is thin.
  \item \textbf{Hidden costs to model:} on-call load, upgrade cadence, expertise ramp, data egress fees.
\end{itemize}

\section{Disaster Recovery Notes (RPO/RTO)}
\begin{itemize}
  \item Define recovery point objective (RPO): maximum acceptable data loss window.
  \item Define recovery time objective (RTO): maximum acceptable restoration time.
  \item Test restores regularly; an untested backup is not a backup.
\end{itemize}

